\documentclass[11pt]{article}
\usepackage[letterpaper,margin=1in]{geometry}
\usepackage[T1]{fontenc}
\usepackage{lmodern,microtype}
\usepackage{amsmath,amssymb,amsthm,mathtools}
\usepackage{xcolor}
\usepackage{fancyhdr}
\usepackage{needspace}
\usepackage[colorlinks=true,linkcolor=blue!42!black,citecolor=blue!42!black,
  urlcolor=blue!42!black]{hyperref}
\hypersetup{
  pdftitle={Twistor spinors, Dirac concentration, and Yamabe spectral rigidity in dimension four},
  pdfauthor={Edward B. Baker III},
  bookmarksdepth=2
}

\pagestyle{fancy}
\fancyhf{}
\fancyhead[L]{\small Twistor spinors and Yamabe spectral rigidity}
\fancyhead[R]{\small Working draft}
\fancyfoot[C]{\thepage}
\setlength{\headheight}{14pt}
\setlength{\emergencystretch}{2em}
\allowdisplaybreaks[1]
\setcounter{tocdepth}{1}
\numberwithin{equation}{section}

\newtheorem{theorem}{Theorem}[section]
\newtheorem{proposition}[theorem]{Proposition}

\DeclareMathOperator{\Ric}{Ric}
\DeclareMathOperator{\Vol}{Vol}
\DeclareMathOperator{\tr}{tr}
\DeclareMathOperator{\ind}{ind}
\DeclareMathOperator{\dist}{dist}
\DeclareMathOperator{\Hess}{Hess}
\DeclareMathOperator{\sn}{sn}
\DeclareMathOperator{\cn}{cn}
\DeclareMathOperator{\dn}{dn}
\newcommand{\rd}{\mathrm{rd}}
\newcommand{\dd}{\mathrm{d}}
\newcommand{\e}{\varepsilon}
\newcommand{\RR}{\mathbb{R}}
\newcommand{\Sphere}{S^4}
\newcommand{\ip}[2]{\langle #1,#2\rangle}
\newcommand{\abs}[1]{\lvert #1\rvert}
\newcommand{\norm}[1]{\lVert #1\rVert}
\newcommand{\dv}{\dd V_g}
\newcommand{\Y}{Y([g])}
\newcommand{\Yf}{Y_f([g])}

\title{Twistor spinors, Dirac concentration, and\\
Yamabe spectral rigidity in dimension four}
\author{Edward B. Baker III}
\date{Working draft --- 25 November 2026}
% Prepared for Edward B. Baker III with OpenAI Codex assistance.
% Actual selected model identifier and reasoning effort were unavailable.
% Revised 24 November 2026 for Edward B. Baker III with LLM assistance:
% Claude Fable 5.2 (claude-fable-5-2), Claude Cowork, reasoning effort max
% as configured. That revision added Section 7.3 on local minimizers and
% conformal vector fields, Section 10 on the flat neck, the Delaunay
% minimizers and the integral Ricci criterion, and the matching changes to
% the abstract, the introduction, Section 9 and the last section.
% Source versions and proof scope are recorded in ../DRAFT_HISTORY.md.

% 25 November 2026 collar integration from checkout 5ff7014, for Edward B. Baker III
% with OpenAI Codex assistance. Actual serving model/effort are unavailable.
% The transmission and comparison statements use assessed MATCH and EXTRACT results.

\begin{document}
\maketitle
\thispagestyle{plain}

\begin{abstract}
We study geometric conditions that recognize the standard smooth
four-sphere. A nontrivial chiral twistor spinor on a closed Riemannian
homotopy four-sphere forces conformal roundness. The least twistor
energy has infimum zero in every conformal class, with concentration of
actual eigenspinors possible. A Dirac Green spinor gives unconditional
boundary data on a punctured homotopy sphere. An exact identity locates
chiral Weyl curvature in a projected derivative of the twistor defect,
and an oscillatory family shows that the first-order error of these
Green fields can tend to zero while the Weyl energy diverges. For a
degree-one map to the round sphere, a modified Yamabe quotient gives a
second variational problem whose unrestricted supremum is approached
by concentrating degree. Imposing moment constraints leads to a
spectral theorem. For a volume-one Yamabe minimizer with scalar
curvature \(s\), the first five positive scalar Laplace eigenvalues
satisfy \(\frac35\sum_{i=1}^5\lambda_i\ge s\), with equality exactly for
a scaled round four-sphere. The fourth-order minimizing inequality
gives a sharp directional moment bound and recovers the metric at
equality. Relative near equality forces the joint law of the five
modes to become spherical, but it does not determine the manifold.
On \(S^1\times S^3\) the Yamabe minimizers are explicit Delaunay
metrics with \(s=3\lambda_1\), and their relative five-eigenvalue gap
tends to zero as the circle lengthens. Attaching a round region across
a conformally flat neck leaves the conformal class unchanged. On a
homotopy four-sphere, a volume-one Yamabe metric whose trace-free Ricci
tensor satisfies \(\int\abs{\Ric_0}^2<s^2/12-16\pi^2\) recognizes the
standard sphere by the theorem of Chang, Gursky, and Yang.
On a supplied two-seam collar, uniform positive Yamabe constant gives
a length-independent Dirac correction for prescribed jumps.
The normalized projected derivative of an actual Green field
measures half the chiral Weyl energy on each measurable source
region. Constructing suitable conformal classes and extracting
controlled comparison geometry on an arbitrary homotopy four-sphere
remain open.
\end{abstract}

\tableofcontents
\clearpage

\section{Introduction}\label{sec:intro}

A geometric approach to the smooth four-dimensional Poincar\'e conjecture
seeks a structure whose existence can be established on an arbitrary
smooth homotopy four-sphere and whose local equations determine its
global smooth type. Freedman's theorem gives the topological
identification~\cite{Freedman1982}. In the smooth Riemannian setting,
the two chiral Weyl blocks give a natural starting point. The twistor
construction of Atiyah, Hitchin, and Singer identifies the vanishing of
one block with integrability of the corresponding complex structure on
the twistor sphere bundle~\cite[Theorem~4.1]{AHS1978}. On a homotopy
sphere the signature identity then removes the other block, and the
conformally flat geometry develops onto the round sphere.

We investigate the first-order spinorial version of this condition.
The ordinary twistor equation recognizes the desired structure through
a global half-spinor. Its variational formulation exposes a loss of
compactness in the metric and spinor. Prescribing a Dirac Green pole
gives a different starting point, with a fixed nonzero boundary value
on a Euclidean end. This formulation identifies the derivative estimate
that a construction would have to supply.

A second use of the Dirac operator couples it to a map to the round
sphere. Its index produces a kernel for every nonzero-degree map, and
scalar comparison turns the geometric condition into a modified
Yamabe problem. Unrestricted maps concentrate their degree in the
same local geometry that gives the sharp quotient. Requiring centered
coordinates brings the first scalar eigenspace into the problem.
Controlling five independent directions leads to a rigidity theorem
for Yamabe minimizers on arbitrary closed four-manifolds, with no spin
assumption. Its stability has a probabilistic conclusion on the
eigenfunction values and a separate geometric requirement for smooth
recognition.

Section~\ref{sec:necklace-tests} tests the construction problem on explicit
cases. Attaching a round sphere to a punctured core across a
conformally flat neck returns the original conformal class, so the
minimizing metrics and their spectra are unchanged. On
\(S^1\times S^3\) the Yamabe minimizers, described by
Schoen~\cite[\S2]{Schoen1989Variational}, are Delaunay metrics that we
write with Jacobian elliptic functions. They satisfy the scalar
equality \(s=3\lambda_1\) at every length, and their relative
five-eigenvalue gap tends to zero on a manifold that is not a sphere.
Finally the Chern--Gauss--Bonnet formula turns the hypothesis of
Chang, Gursky, and Yang into a bound on the trace-free Ricci tensor of
a Yamabe metric, and the Delaunay metrics sit exactly at equality in
that bound.

Section~\ref{sec:collar-transmission} connects the spinor boundary
problem to geometric comparisons of an approximately critical sequence.
When two seams are supplied, the closed Dirac inverse solves the
jump problem in their normalized frames and controls the global
critical spinor norm. A related curvature estimate identifies the
derivative quantity that still needs control on the residual
region, including where the spinor amplitude becomes small.

\section{Chiral recognition and the twistor equation}\label{sec:twistor}

\subsection{The exact geometric condition}

All metrics are smooth and Riemannian unless a limiting regularity is
specified. On an oriented four-manifold let \(\Lambda^\pm\) be the
eigenspaces of the Hodge star on real two-forms. We write \(W^\pm\) for
the blocks of the Weyl operator on these spaces, with \(W^+=0\) the
anti-self-dual condition. Set
\[
 E_\pm(g)=\int_M\abs{W^\pm}_{\mathrm{HS}}^2\,\dv ,
\]
where the norm is the Hilbert--Schmidt norm of the operator in an
orthonormal two-form basis. The full four-index tensor satisfies
\(\abs{W}_{4\mathrm{ind}}^2=4(\abs{W^+}_{\mathrm{HS}}^2+
\abs{W^-}_{\mathrm{HS}}^2)\), since its ordered indices count each
unordered pair twice in each slot. These energies are conformally
invariant in dimension four. In the same convention the signature
formula is
\begin{equation}\label{eq:signature}
 E_+(g)-E_-(g)=12\pi^2\tau(M).
\end{equation}
See Hirzebruch and Gover--Gursky~\cite[\S4]{Hirzebruch1971Signature}
\cite[pp.~293--294]{GoverGursky2024}.

A smooth closed homotopy four-sphere is simply connected, with
\(\tau(M)=0\) and \(\chi(M)=2\). It is spin because \(H^2(M;\mathbb Z/2)=0\),
and its spin structure is unique because \(H^1(M;\mathbb Z/2)=0\).
If \(W^+=0\), \eqref{eq:signature} gives \(W^-=0\). The metric is then
locally conformally flat. In dimension at least three the transition
maps of conformally flat charts are restrictions of M\"obius
transformations. Simple connectivity makes continuation of these charts
a single conformal developing map \(M\to\Sphere\)~\cite[\S1]{SchoenYau1988Conformal}.
It is a local diffeomorphism. Its image is open and, by compactness,
closed in \(\Sphere\). It is therefore a covering, with one sheet
because the target is simply connected. There is also a strict energy
criterion due to Chang, Gursky, and Yang~\cite[Theorem~A]{ChangGurskyYang2003}.
If the conformal class has positive Yamabe constant and
\(E_+(g)<4\pi^2\), then
\(\int\abs{W}_{4\mathrm{ind}}^2\,\dv=8E_+(g)<16\pi^2\chi(M)\).
Their theorem gives a diffeomorphism to \(\Sphere\) or
\(\mathbb{RP}^4\), and simple connectivity selects \(\Sphere\).

\subsection{Spinorial formulation and zeros}

Use \(S^+\) for the complex rank-two half-spin bundle on which the
real self-dual forms act, and \(S^-\) for the opposite chirality.
Write \(c(X)\) for Clifford multiplication, with
\(c(X)^2=-\abs{X}^2 I\), and \(D=\sum_i c(e_i)\nabla_{e_i}\)
for the ordinary spin Dirac operator. The positive twistor operator is
\begin{equation}\label{eq:twistor}
 P_X\psi=\nabla_X\psi+\tfrac14c(X)D\psi .
\end{equation}
Here \(D\) reverses chirality, whereas \(P\) preserves it. If
\(C(T)=\sum_i c(e_i)T_i\), then \(C^*\eta=-\sum_i e^i\otimes c(e_i)\eta\)
and \(CC^*=4I\), so \(P=(I-C^*C/4)\nabla\).
The Schr\"odinger--Lichnerowicz formula and this orthogonal projection
give the identities
\begin{equation}\label{eq:energy}
 \abs{P\psi}^2=\abs{\nabla\psi}^2-\tfrac14\abs{D\psi}^2,\qquad
 P^*P=\tfrac34D^2-\tfrac14s_g,\qquad
 \sigma_2(P^*P)(\zeta)=\tfrac34\abs{\zeta}_g^2 I .
\end{equation}
In particular, ordinary harmonic spinors on closed metrics with
\(s_g>0\) vanish. A pure half-spinor cannot satisfy
\(D\psi=\lambda\psi\) with \(\lambda\ne0\). On the unit round sphere,
Killing spinors satisfy \(\nabla_X\Psi=\kappa c(X)\Psi\) for either
\(\kappa=\pm\tfrac12\)~\cite[\S2]{Bar1996SphereDirac}. Their chiral
projections are nonzero twistor spinors and satisfy
\(D^2\psi=4\psi\). For comparison, the full round Dirac spectrum is
\(\pm(2+k)\), \(k\ge0\). Equation~\eqref{eq:energy} shows that the
higher eigenmodes have strictly positive twistor energy.

Let \(\psi\) be a solution of \eqref{eq:twistor} with \(P\psi=0\).
Set \(\eta=D\psi\) and \(A=\tfrac12(\Ric^\sharp-(s_g/6)I)\).
The prolongation equations are
\begin{equation}\label{eq:prolongation}
 \nabla_X\psi=-\tfrac14c(X)\eta,\qquad
 \nabla_X\eta=-2c(A X)\psi .
\end{equation}
Indeed, taking the covariant trace of the first equation gives
\(\nabla^*\nabla\psi=D\eta/4\), so \(D\eta=s_g\psi/3\).
Commuting the first equation and Clifford-contracting in the second
slot gives
\[
 -\tfrac12 c(\Ric^\sharp X)\psi
   =\sum_i c(e_i)R^S(X,e_i)\psi
   =\tfrac14(2\nabla_X\eta-c(X)D\eta).
\]
The first equality follows by contracting the spin curvature formula
and using the first Bianchi identity. Substitution gives the second
equation in \eqref{eq:prolongation}. In the uncontracted commutator,
its substitution cancels the Schouten part of the spin curvature.
The result is \(W^S(X,Y)\psi=0\), where \(W^S\) denotes the spin
action of the Weyl tensor. The same equations and the zero analysis
appear in Friedrich~\cite[\S2]{Friedrich1990TwistorSpinors}.
At any point where \(\psi\ne0\) a real self-dual form annihilating
\(\psi\) must vanish. Its nonzero action is an invertible traceless
skew-Hermitian matrix on a complex two-plane. Thus \(W^+=0\) on the
nonzero locus.

The pair \((\psi,\eta)\) solves a linear parallel-transport equation
by \eqref{eq:prolongation}. If it vanishes at a point, it vanishes on
the connected manifold. At a zero of a nontrivial \(\psi\),
\(\eta\ne0\) and \(X\mapsto-c(X)\eta/4\) is a real isomorphism from
the tangent space to \(S^+\). The zeros are isolated and \(W^+=0\)
extends across them by continuity.

\begin{proposition}\label{prop:exact-recognition}
On a smooth closed oriented homotopy four-sphere, the existence of a
metric carrying a nontrivial positive chiral twistor spinor is equivalent
to being diffeomorphic to the standard smooth sphere. For a supplied
metric carrying such a spinor, its conformal class is round.
\end{proposition}

\begin{proof}
The forward implication follows from \eqref{eq:prolongation},
the signature formula, and conformal development. For the converse
pull back the round metric and a chiral Killing projection along
an oriented diffeomorphism. Composing with a reflection if necessary
chooses the orientation.
\end{proof}

\subsection{The derivative that contains Weyl curvature}\label{sec:weyl}

There is an exact defect identity for any smooth positive half-spinor.
Let \(T=P\psi\), and use
\[
 m:\Lambda^+\otimes S^+\longrightarrow S^+,\qquad
 m(\omega\otimes v)=c(\omega)v,\qquad \Pi=I-m^*m/6 .
\]
For a real orthonormal self-dual basis \((\omega_a)_{a=1}^3\), put
\(C_a=c(\omega_a)|_{S^+}\). The Clifford relations give
\(C_a C_b+C_b C_a=-4\delta_{ab}I\), so \(mm^*=6I\) and \(\Pi\)
projects to \(\ker m\). For example, with
\(\omega_1=(e^{12}+e^{34})/\sqrt{2}\) the action on \(S^+\) is
\(\sqrt{2}c(e_1)c(e_2)\), and the two other basis elements give the
remaining anticommuting generators.

\begin{proposition}\label{prop:weyl-defect}
With \(\mathcal W(\psi)=\sum_a\omega_a\otimes c(W^+\omega_a)\psi\),
the self-dual projection in the form slot of the exterior covariant
derivative satisfies
\begin{equation}\label{eq:weyl-defect}
 (d^\nabla T)^+=-\tfrac12\mathcal W(\psi)-\tfrac16m^*(P^*T),
 \qquad
 \abs{\Pi(d^\nabla T)^+}^2
   =\tfrac12\abs{W^+}_{\mathrm{HS}}^2\abs{\psi}^2 .
\end{equation}
\end{proposition}

\begin{proof}
Use the positive curvature-operator convention
\(\ip{\mathcal R(X\wedge Y)}{Z\wedge U}=\ip{R(X,Y)U}{Z}\).
The spin curvature equals \(-c(\mathcal R(X\wedge Y))/2\), by the
spin curvature formula~\cite[\S2]{Bar1996SphereDirac}. In a normal frame,
\[
 (d^\nabla T)_{ij}
  =R^S(e_i,e_j)\psi+\tfrac14(c(e_j)\nabla_iD\psi-c(e_i)\nabla_jD\psi).
\]
On contracting the second term with \(\omega_a\), use
\([c(\omega_a),c(e_i)]=2c(\iota_{e_i}\omega_a)\) and
\(c(\omega_a)D\psi=0\) to get \(\tfrac18 C_aD^2\psi\).
The first term contributes \(-\mathcal W(\psi)/2+(s_g/24)m^*\psi\).
The positive curvature block is \(W^++(s_g/12)I\), and the
negative output acts trivially on \(S^+\). Since \(W^+\) is symmetric
and trace-free, the anticommutator implies \(m\mathcal W(\psi)=0\).
Combining the terms with \eqref{eq:energy} proves the first formula.
For a real self-dual \(\omega\) the same algebra gives
\(\abs{c(\omega)\psi}^2=2\abs{\omega}^2\abs{\psi}^2\), so summing
over \(a\) proves the norm identity.
\end{proof}

In particular the adjoint equation \(P^*T=0\) tests the component
orthogonal to the one containing \(W^+\). When the zero set of \(\psi\)
has measure zero, integrating the squared projected derivative divided
by \(\abs{\psi}^2\) gives \(E_+(g)/2\). The integrand has a smooth
extension across the zero set by \eqref{eq:weyl-defect}. On a
homotopy sphere, an integral strictly below \(2\pi^2\) in a
positive Yamabe class is sufficient for smooth recognition.

\section{Variational twistor energy and concentration}\label{sec:twistor-variation}

For a smooth closed spin four-manifold define
\[
 \mu(g)=\inf_{\psi\in H^1(S^+),\ \psi\ne0}
   \frac{\norm{P_g\psi}_{2,g}^2}{\norm{\psi}_{2,g}^2}.
\]
By \eqref{eq:energy} this is the least eigenvalue of a nonnegative
elliptic operator with compact resolvent, attained by a smooth field.
If \(h=a^2g\) and \(\beta:S_g\to S_h\) is the isometry of spin bundles
from orthonormal frame rescaling, then the twistor and Dirac
covariances in four dimensions are
\begin{equation}\label{eq:conformal-spin}
 P_h(a^{1/2}\beta\psi)=a^{-1/2}\mathcal B_a(P_g\psi),\qquad
 D_h(a^{-3/2}\beta\psi)=a^{-5/2}\beta D_g\psi .
\end{equation}
Here \(\mathcal B_a\) is the isometric tensor identification using
the rescaled coframe and \(\beta\). For \(\varphi=\log a\), the
connection on an identified spinor changes in the \(X\)-direction
by \(-c(X)c(\nabla\varphi)/2-X(\varphi)/2\), before replacing the
orthonormal vector by \(a^{-1}X\). Clifford tracing gives the
Dirac weight. Differentiating \(a^{1/2}\beta\psi\) cancels the
terms in the twistor projection. See Hitchin and
Friedrich~\cite[\S1.4]{Hitchin1974Harmonic}\cite[\S1]{Friedrich1990TwistorSpinors}.

\subsection{Zero infimum in every conformal class}

\begin{theorem}\label{thm:twistor-inf}
For every smooth conformal class on a closed connected spin four-manifold,
\[
 \inf_{h\in[g],\ \Vol_h M=1}\mu(h)=0.
\]
The infimum is attained precisely in classes with a nonzero chiral
twistor kernel. In a kernel-free class there are metrics \(h_r\) for
which the mass measures of normalized first eigenspinors converge to
a point mass, and \(\mu(h_r)=O(r^3)\).
\end{theorem}

\begin{proof}
Fix a normal ball at \(p\). Take \(a_r=A_rf_r\), where \(f_r=r^{-1}\)
on \(B_r\), \(f_r=1\) off \(B_{2r}\), and \(1\le f_r\le r^{-1}\).
The factor \(A_r=(\int f_r^4\,\dv)^{-1/4}\) stays between positive
constants by the volume of a small ball. Let \(h_r=a_r^2g\).
A radially parallel unit spinor \(\xi\) satisfies
\(\abs{\nabla^g\xi}=O(\dist_g(\cdot,p))\). In radial gauge the
connection matrix is
\(A_i(x)=\int_0^1 t x^j F_{ji}(tx)\,\dd t\), where \(F\) is its
curvature, giving this bound. Multiply \(\xi\) by a cutoff equal to
one on \(B_{2r}\), supported in \(B_{4r}\), with gradient bounded by
\(C/r\), to obtain \(v_r\). For the field \(a_r^{1/2}\beta v_r\),
the squared numerator and denominator of the Rayleigh quotient
are \(\int a_r^3\abs{P_gv_r}^2\,\dv\) and
\(\int a_r^5\abs{v_r}^2\,\dv\), respectively. The inner ball
contributes \(O(r^3)\) to the numerator and at least \(c/r\) to
the denominator. The cutoff annulus, where \(a_r=A_r\), contributes
\(O(r^2)\) to the numerator. Thus \(\mu(h_r)\le Cr^3\).
Covariance identifies the exact kernels throughout the class.

To check concentration for actual eigenspinors, write a normalized
one as \(a_r^{1/2}\beta u_r\). In a kernel-free class,
\[
 \int_{M\setminus B_{2r}}\abs{a_r^{1/2}\beta u_r}_{h_r}^2\,\dd V_{h_r}
  = A_r^5\int_{M\setminus B_{2r}}\abs{u_r}_g^2\,\dv
  \le C\mu(g)^{-1}\norm{P_gu_r}_{2,g}^2
  \le C'\mu(h_r).
\]
The final bound uses the uniform positive lower bound for \(a_r\).
The measures have mass one and the ball shrinks to \(p\).
\end{proof}

If \(E_+(g)>0\), the family in Theorem~\ref{thm:twistor-inf} has no
positive \(C^2\) limit even after orientation-preserving spin
diffeomorphisms. A limiting metric would have this same Weyl energy
by conformal invariance and \(C^2\) continuity. After the symmetric
frame identifications, the elliptic estimate for \(P\) and
\(C^2\) convergence give a nontrivial weak twistor spinor at the
limit. To see the regularity needed here, \eqref{eq:prolongation}
holds in distributions for a \(C^2\) metric and an \(H^1\) solution,
by approximation in the commutation and contraction identities.
It first gives \(D\psi\in H^1_{\mathrm{loc}}\). The pair then obeys
\(\partial_i U=B_iU\) with continuous local coefficients. Sobolev
embedding improves \(H^1\) first to \(W^{1,4}\) and then to \(W^{1,p}\)
for finite \(p>4\). The continuous pair has continuous first
derivatives by the same equation. The curvature compatibility and
the isolated-zero argument therefore give \(E_+=0\), a contradiction.

\subsection{A conformally invariant relative error}

For a nonzero smooth half-spinor set
\begin{equation}\label{eq:relative}
 J_g(\psi)=\left(\int_{\{\psi\ne0\}}
            \frac{\abs{P_g\psi}^4}{\abs{\psi}^4}\,\dv\right)^{1/4},
\end{equation}
with the empty zero set allowed and the integral permitted to be
infinite. Equation~\eqref{eq:conformal-spin} makes \(J\) invariant
under the twistor transport and under constant rescaling of \(\psi\).
On a fixed kernel-free class, ellipticity, the spectral gap, and
Sobolev embedding give \(\norm{\psi}_4\le C_g\norm{P_g\psi}_2\).
For a field with isolated transverse zeros, H\"older gives
\(\norm{P_g\psi}_2\le J_g(\psi)\norm{\psi}_4\).
Thus \(J\ge C_g^{-1}>0\) throughout that class.

Finite-cost fields exist for every metric on a smooth oriented homotopy
four-sphere. The relation
\(p_1(\Lambda^+)=-4c_2(S^+)=p_1(TM)+2e(TM)\)
follows from the adjoint representation of \(\mathrm{SU}(2)\) and
the sum of the rotation angles of the two tangent planes. It gives
\(\langle c_2(S^+),[M]\rangle=-1\). The real linear model
\(x\mapsto c(x)\eta\), \(0\ne\eta\in S^-_p\), has local index \(-1\).
Indeed \(\mathrm{Spin}(4)\) acts on \(\mathbb H\) by
\(x\mapsto a x b^{-1}\) and may be represented on the half-spinors by
\((s_+,s_-)\mapsto(a s_+,b s_-)\), with Clifford action
\((s_+,s_-)\mapsto(x s_-,-\bar x s_+)\). Complex multiplication on
the right by \(i\) makes the orientation of \(S^+\), represented
by \((1,i,j,-k)\), opposite to the tangent orientation.
Each zero of an exact solution has this index by
\eqref{eq:prolongation}, so the index sum gives exactly one zero.
To extend the local model with only its prescribed zero, remove
a small coordinate ball. Its complement is simply connected and
acyclic by van Kampen and the relative homology sequence, hence
contractible. The local spinor direction in an exterior
trivialization on the boundary has degree zero, since that degree
is the total Euler number minus the local index. Contracting the
direction in a collar, preserving the local data near the seam,
and extending the positive length gives the field. At its zero
\(\abs{\psi}\asymp r\) and \(P\psi(p)=0\) by Clifford trace, so
\(\abs{P\psi}=O(r)\). A transverse zero with incompatible first
jet has \(P\psi(p)\ne0\) and logarithmically divergent cost.

The metric \(\abs{\psi}^{-4}g\) on the punctured manifold carries
the transported unit spinor \(\phi=\abs{\psi}^{-1}\beta\psi\).
Its twistor norm in \(L^4\) equals \(J_g(\psi)\).
For a unit field \(\nabla\phi\) takes values in the real orthogonal
complement of \(\phi\). The Clifford contraction on the complementary
line tensored with covectors is a real isometry to \(S^-\).
Since \(CC^*=4I\), the contraction restricted to the real orthogonal
complement has squared operator norm \(3\). Therefore
\[
 \abs{P\phi}\le\abs{\nabla\phi}\le2\abs{P\phi}.
\]
The quaternionic partner of \(\phi\) completes it to a unit spin
frame. In the matrix norm \(-\tr(A^2)/2\) on the chiral connection,
its connection one-form has norm \(\abs{\nabla\phi}\).
The prescribed zero gives \(\abs{\psi}=\kappa r+O(r^2)\).
In inverted coordinates of radius \(R=r^{-1}\) the metric is
\(\kappa^{-4}\delta+O(R^{-1})\) with derivatives of order \(j\)
bounded by \(O(R^{-1-j})\). In particular the end is complete.
This scale controls the chiral connection in \(L^4\), which controls
its quadratic curvature term in \(L^2\). For an exact solution,
the unit spinor and its quaternionic partner are parallel in this
scale. The vanished positive spin curvature contains the trace-free
Ricci block and \(W^++(s/12)I\), so the punctured metric is Ricci flat.
The derivative in
\eqref{eq:weyl-defect} specifies the separate curvature requirement.

\subsection{A conditional attainment criterion}

Suppose metrics \(g_j\) on a fixed closed spin manifold converge
uniformly and strongly in \(W^{1,4}\) to a smooth positive metric
\(g\), after orientation-preserving spin identifications if needed.
If \(\norm{\psi_j}_{2,g_j}=1\) and either
\(\norm{P_{g_j}\psi_j}_{2,g_j}\to0\) or \(J_{g_j}(\psi_j)\to0\)
with zero sets of measure zero in the relative case, then \(g\)
carries a nontrivial smooth twistor spinor.
Indeed, the symmetric frame comparison identifies the bundles
isometrically~\cite{BourguignonGauduchon1992}. In the fixed bundle
the leading coefficients of \(P\) depend smoothly on the metric and
converge uniformly, and the lower coefficients depend linearly on
its first derivatives with smooth metric-dependent factors.
They converge in \(L^4\). Since \(H^1\) embeds in \(L^4\),
the operators converge in \(H^1\to L^2\). The fixed elliptic estimate
\(\norm{v}_{H^1}\le C(\norm{P_g v}_2+\norm v_2)\) absorbs the difference
for large \(j\). In the relative case one further absorbs
\(C J_{g_j}(\psi_j)\norm{v_j}_{H^1}\) by H\"older and the uniform volume
comparison. The spinors are then bounded in \(H^1\), with a
subsequence strongly convergent in \(L^2\) to a norm-one field.
The weak limiting equation and elliptic regularity give a smooth
solution. Thus the needed metric compactness is independent of the
vanishing scalar residual.

\section{The Dirac boundary problem on a punctured sphere}\label{sec:green}

\subsection{Canonical boundary data and the compact zero}

\begin{proposition}\label{prop:green}
Let \(M\) be a smooth closed oriented homotopy four-sphere and \(p\in M\).
There is a smooth metric \(g\), flat in a ball at \(p\), with invertible
ordinary Dirac operator. For a positive conformal factor \(\omega\) on
\(N=M\setminus\{p\}\) equal to \(r^{-2}\) near \(p\), the metric
\(h=\omega^2g\) has an exactly Euclidean end. Each nonzero constant
positive spinor in an oriented Euclidean end frame is the limit of a
unique harmonic positive spinor with \(O(R^{-3})\) error. Twistor
transport of this field to the compact metric extends smoothly and
has a prescribed transverse zero at \(p\).
\end{proposition}

\begin{proof}
Choose \(g\) flat on a coordinate ball. The complex chiral index is
\(-\tau(M)/8=0\), by the degree-four term of the spin index
formula~\cite{AtiyahSinger1963Index}. In dimension four this gives
vanishing of the index obstruction to invertibility. The relative
metric deformation theorem of Ammann, Dahl, and Humbert gives a metric
with invertible \(D_g\) by changing it in a ball disjoint from the chosen
flat germ~\cite[Theorem~1.1]{AmmannDahlHumbert2011Local}.

Write \(\omega_3=\Vol(S^3)=2\pi^2\) and let
\(\eta\in S^-_p\) be a nonzero source in the parallel spin frame.
The Euclidean fundamental solution of the nonnegative scalar Laplacian
is \(1/(2\omega_3 r^2)\). Applying the flat Dirac operator to this
times \(\eta\) produces its Dirac fundamental solution. Cutting off
that pole and subtracting the inverse Dirac image of the smooth
cutoff error gives the unique compact Green spinor
\begin{equation}\label{eq:green-pole}
 D_gG_\eta=\delta_p\eta,\qquad
 G_\eta(x)=-\frac{c(x)\eta}{\omega_3 r^4}+\gamma(x),
 \qquad D_g\gamma=0\ \text{near }p .
\end{equation}
The remainder is smooth and the field has positive chirality.
This local construction uses only the flat germ and compact
invertibility, as in Ammann, Humbert, and Morel~\cite[\S2]{AmmannHumbertMorel2006Mass}.

By \eqref{eq:conformal-spin}, the field
\(\Psi=\omega^{-3/2}\beta G_\eta\) is \(D_h\)-harmonic. Invert using
\(y=S_a x/r^2\), where \(S_a\) reflects across \(a^\perp\) for a unit vector
\(a\). The reflected inversion preserves orientation, takes the end
to a Euclidean exterior, and has radial coordinate \(R=r^{-1}\).
If \(u=x/r\), its orthonormal frame transition lifts to \(c(a)c(u)\).
In the corresponding end frame
\begin{equation}\label{eq:green-end}
 \Psi_y(y)=\omega_3^{-1}c(a)\eta+
    R^{-3}c(a)c(u)\gamma(S_a y/R^2).
\end{equation}
The map from \(\eta\) to the constant term is an isomorphism. The
error and its derivatives of order \(j\) are \(O(R^{-3-j})\).
If two fields have this asymptotic constant, the inverse harmonic
transport of their difference is bounded near \(p\). Cutting off
a radius-\(\epsilon\) ball when testing its Dirac equation produces
a boundary error of order \(\epsilon^3\). The weak equation holds on
the compact manifold, and elliptic regularity and invertibility give
uniqueness.

The inverse twistor transport uses the different weight in
\eqref{eq:conformal-spin}. It yields
\begin{equation}\label{eq:green-jet}
 \psi=\omega^{-1/2}\beta^{-1}\Psi=\omega^{-2}G_\eta
       =-\omega_3^{-1}c(x)\eta+r^4\gamma(x)
       \quad\text{near }p .
\end{equation}
The leading term is a flat twistor spinor, so the jet is transverse
and \(P_g\psi=O(r^3)\) there.
\end{proof}

We normalize \(\abs{\Psi_\infty}=1\). The complete end has zero ADM
mass and \(\nabla^h\Psi=O(R^{-4})\). Applying the
Schr\"odinger--Lichnerowicz identity to the harmonic field gives
\[
 \tfrac12\Delta_h\abs{\Psi}^2
 =-\abs{\nabla^h\Psi}^2-\tfrac14s_h\abs{\Psi}^2 .
\]
Throughout, \(\Delta=d^*d\) is the nonnegative scalar Laplacian.
Integrate over \(R\le L\). The outward flux of \(\frac12\dd\abs{\Psi}^2\)
is \(O(L^{-1})\), giving
\begin{equation}\label{eq:green-flux}
 \int_N s_h\abs{\Psi}^2\,\dd V_h=-4\int_N\abs{\nabla^h\Psi}^2\,\dd V_h .
\end{equation}
In particular an independently established nonnegative weighted scalar
integral forces parallelism. Equation~\eqref{eq:green-jet} would then
supply an exact compact twistor spinor, including at its zero, and
Proposition~\ref{prop:exact-recognition} would apply. The prescribed
Euclidean end permits this flux calculation without a scalar
curvature sign on the starting conformal class. The positive-mass
construction of Ammann and Humbert uses a scalar-flat end obtained
from a positive Yamabe class~\cite[Theorem~2.2]{AmmannHumbert2005PositiveMass}.

\subsection{Interior zeros and a unit harmonic gauge}

Because \(P_h\Psi=\nabla^h\Psi\), the critical integral on the end
converges. An interior zero makes the same integral infinite.
Indeed, suppose \(f=\abs{\Psi}\) vanishes at the center of a small
coordinate ball. Weak unique continuation for \(D_h^2\) gives a
nearby ball where \(f\) is positive~\cite{Aronszajn1957Continuation}.
If the critical integral were finite, the functions
\(u_L=\log(\max(f,e^{-L}))\) would have uniformly bounded
\(L^4\) gradient by \(\abs{\dd f}\le\abs{\nabla\Psi}\).
Poincar\'e's inequality and the nearby ball on which their values are
bounded give a uniform \(W^{1,4}\) bound. Fatou and weak compactness
give \(u=\log f\in W^{1,4}\) almost everywhere, including that the
zero set has measure zero. If \(\bar u(r)\) is the spherical mean,
H\"older on the radial integral gives
\[
 \abs{\bar u(r)-\bar u(r_0)}
 \le C\norm{\dd u}_{L^4(B)}\bigl(\log(r_0/r)\bigr)^{3/4}.
\]
Smoothness of \(\Psi\) gives \(f\le Cr\) and hence
\(\bar u(r)\le \log(Cr)\), contradicting this estimate.
Conversely, a zero-free field has finite critical integral by
smoothness, positivity on the compact part and the end decay.
By twistor covariance this is also the criterion for finite
\(J_g(\psi)\) for the compact field \eqref{eq:green-jet}.

Each half-spin fiber is a quaternionic line, and the Dirac Green
operator intertwines the invariant quaternionic structures. Evaluation
of \(G_\eta\) at a point off the pole is therefore a quaternionic-linear
map of lines. Changing the unit source multiplies the result on
the right by a unit quaternion. Its norm, zero set and relative
energy are unchanged. If the interior zeros are transverse, their
signed sum is zero because the pole in \eqref{eq:green-jet} already
accounts for \(\langle c_2(S^+),[M]\rangle=-1\).
Both signs are possible for local ordinary harmonic spinors.
For example, in quaternionic tangent axes
\(e_\mu=(1,i,j,k)\), a diagonal linear field
\(\sum_\mu d_\mu x_\mu e_\mu\) is harmonic when
\(\sum_\mu d_\mu=0\). The choices \((1,1,1,-3)\) and
\((1,1,-1,-1)\) have opposite transverse indices.

For a zero-free field put \(\sigma=\abs{\Psi}_h\) and
\[
 k=\sigma^{4/3}h,\qquad \xi=\sigma^{-1}\beta_{h,k}\Psi,\qquad
 V=\abs{\nabla^k\xi}^2 .
\]
The harmonic weight makes \(\xi\) a unit \(D_k\)-harmonic spinor,
and the norm identity gives \(s_k=-4V\).
In the end coordinates \(k_{ij}=\delta_{ij}+O_2(R^{-3})\) and
\(V=O(R^{-8})\), with \(O_\ell\) indicating corresponding derivative
decay through order \(\ell\). For \(v>0\) equal to one off a compact
set take \(k_v=v^{-4}k\) and \(\xi_v=v^3\beta_{k,k_v}\xi\).
The energy and the critical cost are then
\begin{equation}\label{eq:harmonic-energy}
 \int_N\abs{\nabla^{k_v}\xi_v}^2\,\dd V_{k_v}
 = Q(v):=\int_N (Vv^2+12\abs{\dd v}^2)\,\dd V_k,\qquad
 J_{k_v}(\xi_v)^4=\int_N (V+12\abs{\dd\log v}^2)^2\,\dd V_k .
\end{equation}
For the coefficient \(12\), put \(A=\nabla^k\log v\).
In the rescaled orthonormal frame the differentiated spinor,
after removing the factor \(v^5\), has components
\(\nabla^k_i\xi+c(e_i)c(A)\xi+4A_i\xi\).
The cross term with the first component sums to zero by \(D_k\xi=0\)
and the unit norm. The square of the correction sums to
\((4+16-8)\abs A^2\), since
\(\operatorname{Re}\ip{c(e_i)c(A)\xi}{\xi}=-A_i\).
Using \(\dd V_{k_v}=v^{-8}\dd V_k\) and \(\abs{\xi_v}=v^3\) gives
both formulas.

\subsection{Fixed-geometry optimization and loss of core scale}

For fixed \(k\) and \(V\) as above, the infimum \(P_0\) of \(Q(v)\)
over the positive compact variations in \eqref{eq:harmonic-energy}
is positive if \(V\) is nonzero. Let \(\dot H^1(N)\) be the closure
of compactly supported smooth functions in the gradient norm. The
homogeneous inequality \(\norm{w}_4\le S\norm{\dd w}_2\) holds on
\(N\). To see this, cut off to the end and use the Euclidean Sobolev
inequality. The cutoff error in the fixed transition annulus is
controlled by the gradient via integration to infinity on rays
and Cauchy--Schwarz with weight \(R^3\). A compact Poincar\'e
inequality with the mean anchored on that annulus controls the
remaining compact contribution. If \(V\ge c>0\) on a ball of
volume \(b>0\), then every \(v\in1+\dot H^1(N)\) satisfies
\[
 b^{1/2}\le b^{1/4}S\norm{\dd v}_2+
    c^{-1/2}\Big(\int Vv^2\,\dd V_k\Big)^{1/2}.
\]
This proves the positive lower bound.

The bilinear form \(12\int\ip{\dd w}{\dd z}+\int Vwz\) is continuous
and coercive on \(\dot H^1\). Solving there with the linear functional
\(\int Vz\) gives a unique smooth relaxed minimizer \(v_*=1-w\)
with
\begin{equation}\label{eq:cap-eqn}
 12\Delta_k v_*+Vv_*=0,\qquad
 P_0=Q(v_*)=\int_N Vv_*\,\dd V_k .
\end{equation}
Truncation to \([0,1]\) reduces the form and the strong maximum
principle gives \(0<v_*\le1\), with \(v_*<1\) everywhere if \(V\ne0\).
Smooth compact truncations of \(w\) approximate in \(\dot H^1\)
by the \(L^4\) control of the shell cutoff error. The corresponding
positive \(v\)'s converge in energy. The identity for \(P_0\) follows
by testing with \(1-v_*\). Thus the relaxed minimum equals the
original infimum, and the original class attains it exactly for \(V=0\).

The asymptotic expansion and its first derivative are
\begin{equation}\label{eq:cap-mass}
 v_*=1-\frac{P_0}{24\omega_3}R^{-2}+O_1(R^{-3}),\qquad
 m_{\mathrm{ADM}}(v_*^{-4}k)=\frac{P_0}{6\omega_3}.
\end{equation}
For the expansion, \(w=1-v_*\) solves
\(\Delta_k w=Vv_*/12=O(R^{-8})\). A multiple of \(R^{-2}-R^{-3}\)
is an exterior supersolution at large \(R\). The weak maximum
principle in \(\dot H^1\) gives \(w=O(R^{-2})\), and interior
elliptic estimates on rescaled annuli give \(\dd w=O(R^{-3})\).
After an exterior cutoff to \(\RR^4\), the Euclidean equation has
the form \(\Delta_0 w=F+\partial_i G_i\) with
\(F=O(R^{-8})\), \(G=O(R^{-6})\). The flux of
\(\dd v_*\) is \(P_0/12\) by \eqref{eq:cap-eqn}, so
\(\int F=P_0/12\). Convolution with \(1/(2\omega_3 R^2)\)
gives \(P_0/(24\omega_3 R^2)+O(R^{-3})\). The scalar source has
finite first moment and the divergence source contributes
\(O(R^{-3})\). These estimates follow by splitting each convolution
at half the observation radius and using the stated decay on the
tail. The difference from this convolution is harmonic and decays
at infinity. Rescaled elliptic estimates for the remainder give
the first derivative decay, and annular \(W^{2,p}\) bounds follow
from the same equation. With the normalization
\[
 m_{\mathrm{ADM}}(\gamma)
 =\frac{1}{6\omega_3}\lim_{L\to\infty}
 \int_{S_L}(\partial_j\gamma_{ij}-\partial_i\gamma_{jj})
              \nu^i\,\dd S ,
\]
the scalar coefficient of \(R^{-2}\delta\) in \(v_*^{-4}k\)
is \(P_0/(6\omega_3)\) and equals its mass. The decay bounds and
the scalar transformation law give integrable scalar curvature
for this metric, as required for the usual coordinate invariance
of this flux~\cite[\S4]{Bartnik1986Mass}. In fact
\(s_{k_{v_*}}=-3Vv_*^4-36v_*^2\abs{\dd v_*}^2\).
Compact variations keep zero mass. At the relaxed minimizer the
spinor boundary term in the integrated identity is
\(3\lim_{L\to\infty}\int_{S_L}v_*\partial_\nu v_*=P_0/4\).
This agrees with Dirichlet energy \(P_0\) and weighted scalar
integral \(-3P_0\).

Allowing the metric to vary destroys a uniform positive lower bound
for this Dirichlet energy even at fixed end normalization. For
\(0<t\le1\), choose an orientation-preserving radial diffeomorphism
\(\Phi_t\) equal to the identity on the core and to \(y\mapsto y/t\)
far out in the Euclidean end. A nondecreasing radial interpolation
gives a positive radial derivative for all \(t\). Take
\(h_t=t^2\Phi_t^*h\) and transport \(\Phi_t^*\Psi\) by the unitary
spin identification for the constant scale. Then \(h_t\) is exactly
Euclidean in the same exterior chart and
\(\Psi_t(y)=\Psi(y/t)\) there, with the same constant term. In the
compact chart \(\Phi_t\) extends smoothly across \(p\) as \(x\mapsto tx\).
The metric \(g_t=\omega^{-2}h_t\) has the same flat germ and is
conformal to \(\Phi_t^*g\). The compact Green field has the same
singular term in \eqref{eq:green-pole} and remainder \(t^3\gamma(tx)\)
near the pole, by uniqueness. Change of variables gives
\[
 \int\abs{\nabla^{h_t}\Psi_t}^2\,\dd V_{h_t}
   =t^2\int\abs{\nabla^h\Psi}^2\,\dd V_h,\qquad
 \int \frac{\abs{\nabla^{h_t}\Psi_t}^4}{\abs{\Psi_t}^4}\,\dd V_{h_t}
   =\int \frac{\abs{\nabla^h\Psi}^4}{\abs{\Psi}^4}\,\dd V_h ,
\]
where the second integral is over the respective nonzero loci.
The core volume and its spinor \(L^2\) mass both scale by \(t^4\).
Thus the unrestricted Dirichlet infimum is zero on every input,
and zero-energy attainment would furnish an exact compact solution.

\subsection{An oscillatory obstruction for the actual Green fields}

\begin{proposition}\label{prop:oscillation}
On the standard smooth \(\Sphere\) there are compact metrics \(g_n\)
with a common flat pole germ and invertible Dirac operator whose
normalized positive Green fields have no interior zeros. Their compact
twistor transports \(\psi_n\) have the same transverse pole jet and
\[
 g_n\longrightarrow g_b\quad\text{in }C^{1,1/4},\qquad
 J_{g_n}(\psi_n)=O(n^{-1/2}),\qquad E_+(g_n)\ge c n .
\]
The scalar curvature in the corresponding unit harmonic scales
has \(L^2\) norm \(O(n^{-1})\).
\end{proposition}

\begin{proof}
Use \(N\cong\RR^4\) and choose the compact metric \(g_b=\omega^{-2}\delta\),
where \(\omega\) equals \(1\) on a bounded open patch \(U\) and
equals the square of the Euclidean radius at infinity.
Then \(g_b\) is conformal to round and smooth and flat near the
pole. Its Dirac operator is invertible by conformal covariance of
the round kernel, and the punctured Green field is the constant unit
spinor \(\kappa\).

In \(U\) take \(H=(\dd x^2)^2-(\dd x^3)^2\),
\(T=(\dd x^1)^2\circledast H\), and \(C=2T-H\circledast\delta\),
where \(\circledast\) is the Kulkarni--Nomizu product and
\(T_{1a1b}=H_{ab}\) for \(a,b>1\). The tensor \(C\) is
algebraic Weyl, with \(C_{i1j1}=H_{ij}\). For a nonzero smooth
real cutoff \(\chi\) supported in \(U\), set
\[
 F_n=\chi\cos(nx^1),\qquad
 (b_n)_{ij}=-n^{-7/2}\sum_{k,l}C_{ikjl}\partial_k\partial_l F_n,
 \qquad g_n=g_b+b_n .
\]
The algebraic symmetries make \(b_n\) trace-free and divergence-free
on the Euclidean patch. Its leading term is \(n^{-3/2}F_n H\)
with a remainder \(O(n^{-5/2})\) whose \(j\)-th derivatives are
\(O(n^{j-5/2})\) for \(j\le2\). In particular \(g_n\to g_b\) in
\(C^{1,1/4}\) by interpolation of the first and second derivative
bounds. These metrics are positive for large \(n\).

Identify the spinors by the symmetric frame maps. The difference of
the identified \(D_n\) from \(D_b\) is \(O(n^{-1/4})\) from
\(C^{1,1/4}\) to \(C^{0,1/4}\) and \(O(n^{-1/2})\) from \(W^{1,4}\)
to \(L^4\), by the first-order coefficient bounds. The inverse of
\(D_b\) on both scales is bounded by elliptic regularity. Neumann's
series gives invertibility and uniform bounds for the two inverses
of \(D_n\). Use the same cutoff parametrix \(Q\) for the pole.
If \(G_b-Q=S_b\) is smooth and \(G_n-G_b=u_n\) in the identified
bundle, then
\[
 D_n u_n=-(D_n-D_b)S_b,\qquad
 \norm{u_n}_{C^{1,1/4}}=O(n^{-1/4}),\qquad
 \norm{u_n}_{W^{1,4}}=O(n^{-1/2}).
\]
On a smaller pole ball \(u_n\) is flat Dirac-harmonic, and interior
estimates there give fixed-order derivatives \(O(n^{-1/2})\).
It follows from \eqref{eq:green-end} and the compact estimates that
the end fields \(\Psi_n\) converge uniformly to \(\kappa\).
The connection difference on the perturbation patch is
\(O(n^{-1/2})\) and the end derivative has bound
\(O(n^{-1/2}R^{-4})\). Thus
\(\norm{\nabla\Psi_n}_{L^2}+\norm{\nabla\Psi_n}_{L^4}=O(n^{-1/2})\).
Uniform nonvanishing and twistor covariance give the asserted bound
for \(J\). The unit harmonic scales have the same \(L^4\) bound on
their spinor derivatives, by the conformal connection law and
\(\abs{\dd\abs{\Psi_n}}\le\abs{\nabla\Psi_n}\).
The relation \(s_k=-4V\) gives the scalar bound.

In the curvature calculation, both differentiations falling on the
leading cosine give the linear contribution
\(\frac12 n^{1/2}F_n T\), by the flat curvature formula
\(R'_{1a1b}=-\partial_1^2(b_n)_{ab}/2\) at this order.
Its Weyl projection is \(\frac14n^{1/2}F_n C\).
All cutoff terms and differentiated remainders are \(O(n^{-1/2})\).
The nonlinear curvature error is
\(O(\abs{b_n}\abs{\partial^2b_n}+\abs{\partial b_n}^2)=O(n^{-1})\).
Since \(\int F_n^2\to\frac12\int\chi^2>0\), the full four-index
Weyl energy is bounded below by a positive multiple of \(n\).
Equation~\eqref{eq:signature} and the norm convention give \(E_+\)
as one eighth of that energy. This proves the claimed failure
of first-order control even for the canonical boundary fields.
\end{proof}

\section{Twisted Dirac comparison and the modified Yamabe quotient}\label{sec:twisted}

Let \((M,g)\) be a closed connected oriented spin four-manifold.
Let \(h_{\rd}\) be the unit round metric on \(\Sphere\), and let
\(f:M\to\Sphere\) have nonzero degree. Such maps of degree one exist
on every oriented \(M\) by collapsing the complement of an oriented
coordinate ball, using a smooth representative constant on a collar.
Write
\begin{align}
 V_f&=3\abs{df}_{g,h_{\rd}}^2,\qquad
 L_{g,f}=6\Delta_g+s_g-V_f,\notag\\
 Q_f(u)&=\int_M \bigl(6\abs{\dd u}^2+(s_g-V_f)u^2\bigr)\,\dv,\qquad
 \Yf=\inf_{u\in H^1(M)\setminus\{0\}}\frac{Q_f(u)}{\norm{u}_{4,g}^2}.
 \label{eq:modified}
\end{align}
All differential norms of maps are Hilbert--Schmidt norms.
The ordinary Yamabe constant \(\Y\) uses the same quotient with
\(V_f=0\). The conformal scalar formula in dimension four gives,
for \(k=a^2g\),
\begin{equation}\label{eq:conformal-L}
 L_{k,f}(u)=a^{-3}L_{g,f}(au),\qquad
 s_k-V_{k,f}=a^{-3}L_{g,f}a .
\end{equation}
The form and \(L^4\) norm transform by multiplication of the test
function by \(a\). In particular \(\Yf\) depends only on the class
and the fixed map. See Lee--Parker and Listing for the Yamabe
covariance and the fixed-potential functional~\cite[\S2]{LeeParker1987Yamabe}
\cite[\S2.2]{Listing2012ScalarBundles}.

\subsection{Index, curvature inequality, and spectral sign}

Twist the domain spin connection by the pullback of the spin
connection on \(S_{h_{\rd}}=S_{h_{\rd}}^+\oplus S_{h_{\rd}}^-\).
For \(E^\pm=f^*S_{h_{\rd}}^\pm\), the degree-four term of the target
Chern character is \((p_1\pm 2e)/4\). This follows by squaring the
formal half-spin weights \(\pm(x_1\pm x_2)/2\).
On the sphere \(p_1=0\) by stable triviality and \(\langle e,[S^4]\rangle=2\).
The product-graded index is consequently
\[
 \ind D^{E^+,+}-\ind D^{E^-,+}=2\deg f
\]
by the spin index formula~\cite{AtiyahSinger1963Index}. On a
homotopy sphere the untwisted chiral index is zero, and the two
individual indices are \(\deg f\) and \(-\deg f\), respectively.
The product index gives a nonzero twisted harmonic spinor for
every such map and every supplied metric.

We use the sphere comparison of Listing~\cite[Theorem~3]{Listing2010ScalarSymmetric}.
Here is a direct proof for the scalar bound \(s_g\ge V_f\).
In singular frames let \(\sigma_1,\ldots,\sigma_4\) be the singular
values of \(df\). The twist curvature is
\(\sum_{i<j}c(e_i)c(e_j)\otimes R^E_{ij}\). Each summand has operator
norm \(\sigma_i\sigma_j/2\) by round curvature one. Thus the
Schr\"odinger--Lichnerowicz formula for a twisted harmonic spinor
\(\Phi\) gives
\[
 0\ge \int_M\left(\abs{\nabla\Phi}^2+
              \frac{s_g-b}{4}\abs{\Phi}^2\right)\,\dv,\qquad
 b=\sqrt{24\sum_{i<j}\sigma_i^2\sigma_j^2}\le 3\sum_i\sigma_i^2.
\]
The first bound uses Cauchy--Schwarz on the six summands.
The second follows from
\(\sum_{i<j}a_i a_j\le 3(\sum_i a_i)^2/8\) for \(a_i=\sigma_i^2\).
If \(s_g\ge V_f\ge b\), the spinor is parallel with positive constant
norm and all scalar bounds are equalities. Hence the singular
values coincide and \(f^*h_{\rd}=\alpha g\), with
\(s_g=12\alpha\) for a smooth nonnegative \(\alpha\).
Where \(\alpha>0\), the pulled back metric is round locally, so the
scalar formula gives \(\Delta_g\sqrt{\alpha}=0\).
Equivalently \(2\alpha\Delta_g\alpha+\abs{\dd\alpha}^2=0\) there.
The latter identity holds on \(\{\alpha=0\}\) as well because
\(\dd\alpha=0\) at a zero. Integrating gives
\(3\int_M\abs{\dd\alpha}^2\,\dv=0\). The degree ensures
\(\alpha>0\) somewhere. Thus \(\alpha\) is positive and constant
and the map is a homothetic covering of the simply connected sphere.

Let \(\ell_f(g)\) be the least eigenvalue of \(L_{g,f}\).
For its positive ground state \(a\), \eqref{eq:conformal-L} gives
\(s_{a^2g}-V_{a^2g,f}=\ell_f(g)a^{-2}\).
If \(\ell_f(g)\ge0\), the preceding rigidity forces a homothety in
the new representative and the deficit vanishes. Conversely, when
\(g\) is conformal to the pullback by a diffeomorphism, the constant
ground state for the round pullback transports to a positive
zero eigenfunction. Consequently
\begin{equation}\label{eq:sign}
 \ell_f(g)\le0,\qquad
 \ell_f(g)=0\quad\Longleftrightarrow\quad
 f\text{ is a conformal diffeomorphism}.
\end{equation}
A ground-state test also gives \(\Yf\le0\), with equality
precisely in the equality case of \eqref{eq:sign}.

\subsection{Sharp optimization with the map fixed}

\begin{proposition}\label{prop:fixed-map}
For a fixed nonzero-degree \(f\),
\[
 \max_{\widetilde g\in[g],\ \Vol_{\widetilde g}M=1}
       \ell_f(\widetilde g)=\Yf .
\]
If \(\Yf<0\), there is a unique maximizing representative, and
it has constant scalar deficit \(s_{\widetilde g}-V_{\widetilde g,f}=\Yf\).
If \(\Yf=0\), every volume-one representative maximizes the
eigenvalue, and the normalized round pullback is the unique
representative with identically zero pointwise deficit.
\end{proposition}

\begin{proof}
For \(Y_f<0\), a nonnegative minimizing sequence with
\(\norm{u_j}_4=1\) is bounded in \(H^1\) by boundedness of the potential
and compactness of \(M\). Pass to a weak \(H^1\) and strong \(L^2\)
limit \(u\). Lower semicontinuity gives
\(Y_f\le Y_f\norm u_4^2\le Q_f(u)\le Y_f<0\), so the full critical
norm survives. The Euler equation is \(L_{g,f}u=Y_f u^3\).
Its right-hand cubic term is nonpositive. In particular \(u\)
is a nonnegative subsolution of \((6\Delta_g+1)u\le Cu\) for
a constant \(C\). The solution of \((6\Delta_g+1)v=Cu\)
belongs to \(W^{2,4}\) and is bounded. Testing with \((u-v)_+\)
gives \(u\le v\). Thus elliptic bootstrap and the strong maximum
principle give a smooth positive \(u\). For \(Y_f=0\), use a
positive ground state normalized in \(L^4\).

In both cases \(u^2g\) has volume one and constant deficit \(Y_f\)
by \eqref{eq:conformal-L}. For \(a^2g\) of volume one, test the
ordinary eigenvalue quotient with \(u/a\) to obtain
\[
 \ell_f(a^2g)\le\frac{Y_f}{\int_M a^2u^2\,\dv}\le Y_f ,
 \qquad 0<\int_Ma^2u^2\,\dv\le1 .
\]
For \(Y_f<0\), equality in Cauchy--Schwarz gives \(a=u\).
For \(Y_f=0\), the spectral equality is conformally invariant by
\eqref{eq:sign}. A representative with zero pointwise deficit is a
homothetic pullback by the fixed \(f\), and its scale is fixed by
the volume.
\end{proof}

\subsection{The sharp unrestricted supremum over degree-one maps}

\begin{theorem}\label{thm:map-sup}
For any conformal class on a closed connected oriented spin four-manifold,
\[
 \sup_{\deg f=1}\Yf=\min\{\Y,0\}.
\]
Attainment is equivalent to the class being conformally round. The
maximizers are the orientation-preserving conformal diffeomorphisms.
There is a maximizing sequence whose signed Jacobian measures
converge to \(\Vol(S^4,h_{\rd})\delta_p\), for any fixed \(p\).
\end{theorem}

\begin{proof}
The upper bound follows from \eqref{eq:sign} and \(Q_f\le Q_0\),
where \(Q_0\) denotes the ordinary form without the map.
For the lower bound, use inverse stereographic projection in
normal coordinates at \(p\), with
\[
 F_\e(x)=\left(\frac{2\e x}{\e^2+r^2},
               \frac{\e^2-r^2}{\e^2+r^2}\right),\qquad
 b_\e(x)=\frac{2\e}{\e^2+r^2}.
\]
It has conformal factor \(b_\e\) in the Euclidean chart. For
\(\rho=\e^{1/2}\) and \(R=\e^{1/4}\), use the exact map inside
\(B_\rho\). Cut it off on \(\rho<r<2\rho\) in the south-pole coordinates
by multiplying \(\e x/r^2\) by a radial cutoff, and continue by
the south pole. The cutoff map \(f_\e\) has degree one by the north
pole preimage, with the orientation chosen there. Its derivative
on the cutoff annulus is \(O(\e/\rho^2)\).
The model potential is \(p_\e=12b_\e^2\), and
\(6\Delta_0b_\e=p_\e b_\e\). The ground-state identity gives
\begin{equation}\label{eq:bubble-form}
 \int_{\RR^4}(6\abs{\dd v}^2-p_\e v^2)\,\dd x
   =6\int_{\RR^4}b_\e^2\abs{\dd(v/b_\e)}^2\,\dd x\ge0
\end{equation}
for compactly supported tests.

The model potential has bounded \(L^2\) norm, with
\(L^2\) norm \(O((\e/\rho)^2)\) outside \(B_\rho\).
The annular derivative bound gives the same error for the actual
potential there. In \(B_\rho\) the normal-coordinate metric and
volume introduce a factor \(1+O(R^2)\). Thus on \(B_R\) the
actual potential times the volume density differs from \(p_\e\)
in \(L^2\) by \(O(R^2+(\e/\rho)^2)\). The scalar-curvature term
has \(L^2\) norm \(O(R^2)\). In comparing the local kinetic energy
keep the factor \(1-CR^2\) on the full nonnegative form in
\eqref{eq:bubble-form}. H\"older bounds the remaining potential
error uniformly in the test gradient. Consequently, for
\(v\in H^1_0(B_R)\),
\begin{equation}\label{eq:local-bubble}
 Q_{f_\e}(v)\ge -C\bigl(R^2+(\e/\rho)^2\bigr)\norm v_{4,g}^2 .
\end{equation}

Write a partition of squared cutoffs as
\(\chi^2+\theta^2=1\), where \(\chi=1\) for \(r\le2\rho\)
and \(\chi=0\) for \(r\ge R\). Taking the sine and cosine of a
cutoff varying on the logarithmic radial interval gives
\(\norm{\dd\chi}_4^4+\norm{\dd\theta}_4^4=O(\log(R/(2\rho))^{-3})\).
The gradient cross terms cancel because
\(\chi\,\dd\chi+\theta\,\dd\theta=0\). The localized exterior
test has potential zero, so
\[
 Q_{f_\e}(u)=Q_{f_\e}(\chi u)+Q_0(\theta u)
     -6\int_M(\abs{\dd\chi}^2+\abs{\dd\theta}^2)u^2\,\dv
 \ge \Y\norm{\theta u}_4^2-o(1)\norm u_4^2 .
\]
For \(\Y\ge0\) the exterior term is nonnegative. For \(\Y<0\) it
is bounded below by \(\Y\norm u_4^2\). This proves the limit
uniformly over the \(L^4\) unit sphere. The inner Jacobian is
\(b_\e^4\,\dd x\), whose integral tends to the sphere volume by
\(x=\e z\). The annulus contributes \(O((\e/\rho)^4)\), so
the degree measure concentrates at \(p\).
The same scaling gives
\(V_{f_\e}^2\,\dv\rightharpoonup144\Vol(S^4,h_{\rd})\delta_p\).

If \(\Y<0\), test \(Q_f\) with a positive ordinary Yamabe minimizer.
Since \(f\) has nonzero degree, the subtraction is strict, so
\(\Yf<\Y\). Existence and positivity of the ordinary Yamabe
minimizer follow from the Yamabe theorem~\cite{LeeParker1987Yamabe,Schoen1984Yamabe}.
If \(\Y\ge0\), attainment means \(\Yf=0\), hence \eqref{eq:sign}
gives a conformal diffeomorphism. Such a map realizes zero
and its class has positive Yamabe constant.
\end{proof}

\section{Balanced maps, spectral limits, and degree}\label{sec:balanced}

Fix a smooth volume-one Yamabe minimizer \(g\), so its constant
scalar curvature is \(s=\Y\). Unless specified otherwise, all
integrals in the fixed-metric arguments use \(\dv\), and
\(\lambda_i\) denotes the \(i\)-th positive eigenvalue of
\(\Delta_g=d^*d\), counted with multiplicity. We first assume
\(M\) oriented and spin for the degree comparison. A map
\(f:M\to S^4\subset\RR^5\) is called balanced when
\(\int_M f\,\dv=0\). Balanced degree-one maps exist for every
smooth positive volume form after normalization, as shown in
Section~\ref{sec:moment-normalization}.

For an \(L^2\)-normalized first eigenfunction \(\phi\), the Yamabe
inequality applied to \(1+t\phi\) gives \(6\lambda_1-2s\ge0\)
at order \(t^2\). In fact the numerator is
\(s+(6\lambda_1+s)t^2\), and the squared \(L^4\) norm is
\(1+3t^2+O(t^3)\). It follows that
\begin{equation}\label{eq:balance-chain}
 \Yf\le s-3\int_M\abs{df}^2\,\dv\le s-3\lambda_1\le0
 \quad\text{for balanced } f .
\end{equation}
The first comparison uses the constant test and the second
uses Poincar\'e on each coordinate and \(\abs f=1\).
If \(3\lambda_1>s\), this gives a strict uniform gap on the
balanced class.

\subsection{Characterization at zero balanced supremum}

\begin{theorem}\label{thm:balance}
On the supplied minimizing metric, the supremum of \(\Yf\)
over balanced degree-one maps is zero if and only if
\(s=3\lambda_1\) and there exists a smooth
\(F:M\to S^4\) satisfying \(\Delta_g F=\lambda_1 F\)
componentwise. Every balanced degree-one maximizing sequence at
zero has a subsequence converging strongly in ambient \(W^{1,2}\)
to such an eigenmap. A nonzero-degree limit is a homothetic
diffeomorphism. If the absolute Jacobians of the sequence are
uniformly integrable, the limit has degree one.
\end{theorem}

\begin{proof}
When \(\Yf\to0\), \eqref{eq:balance-chain} gives
\(s=3\lambda_1>0\) and \(\int\abs{df}^2\to\lambda_1\).
On the orthogonal complement of the constants and first
eigenspace, the spectral gap bounds the \(H^1\) norm squared
of that part of \(f\) by a fixed multiple of
\(\int\abs{df}^2-\lambda_1\). The projections onto the first
eigenspace admit a convergent subsequence because it is finite
dimensional. Thus the limit \(F\) is smooth, has \(\abs F=1\),
and satisfies the stated eigenvalue equation. The product rule
for \(\abs F^2\) gives \(\abs{dF}^2=\lambda_1\). Therefore
\(s=3\abs{dF}^2\), and the scalar comparison in
Section~\ref{sec:twisted} gives homothety when \(\deg F\ne0\).
The determinant formula for the signed Jacobian is continuous
in \(F,dF\), so the Jacobians converge in measure.
Uniform integrability upgrades this to \(L^1\) convergence,
preserving the degree. For example, a uniform \(L^{4+\eta}\)
derivative bound for any \(\eta>0\) suffices by
\(\abs{J_f}\le\abs{df}^4/16\).

For the converse, a supplied first eigenmap at \(s=3\lambda_1\)
is balanced and has pointwise zero modified deficit. Its degree
is \(0,1\), or \(-1\) by the scalar comparison. If the degree
is one, it attains zero. Otherwise set \(k=1-\deg F\) and
choose \(k\) disjoint coordinate balls.
In target normal coordinates flatten \(F\) to \(F(p_i)\) through
radius \(3\rho\) around each center, restoring \(F\) at \(4\rho\).
This is a homotopy with bounded derivative on the flattened
annulus, since in the target chart the original coordinate
distance is \(O(\rho)\). Insert the degree-one bubbles of
Theorem~\ref{thm:map-sup}, rotated to have the same constant
on the insertion collars. The resulting maps \(G_\e\) have
degree one. Each bubble has square derivative integral
\(O(\e^2(1+\log(\rho/\e)))\), by integration of \(b_\e^2\),
so \(G_\e\to F\) strongly in \(W^{1,2}\).
Their mean is \(O(\rho^4)\).

On each localization ball \(B_R(p_i)\) the bubble-potential
estimate leading to \eqref{eq:local-bubble} still holds. The
bounded derivatives on the outer part contribute \(O(R^2)\)
in the \(L^2\) potential comparison. Off the insertion and
flattening balls the unmodified deficit is zero. A squared
partition with the same logarithmic cutoffs therefore gives
\(Q_{G_\e}(u)\ge-o(1)\norm u_4^2\) uniformly, and
\(\norm{V_{G_\e}}_2\) is bounded.

To impose exact balance, use the time-one conformal flow
\(\Phi_b\) of \(b-(b\cdot y)y\) on \(S^4\) for \(b\in\RR^5\)
near zero. The field is the spherical gradient of \(b\cdot y\),
whose Hessian is \(-(b\cdot y)h_{\rd}\).
The derivative at \(b=0\) of the mean map
\(b\mapsto\int\Phi_b(F)\,\dv\) is
\(I-\int FF^{\mathsf T}\,\dv\). It is positive definite because
equality in a direction would confine the connected image of
the nonconstant \(F\) to two antipodal points.
The mean maps for \(G_\e\) converge in \(C^1\) on a small parameter
ball and their mean errors at zero are \(O(\rho^4)\). Quantitative
local inversion thus gives \(b_\e=O(\rho^4)\) with
\(\int\Phi_{b_\e}(G_\e)\,\dv=0\).
The target conformal factor is \(1+O(\abs{b_\e})\) in \(C^0\),
so the potential changes by \(o(1)\) in \(L^2\). This preserves
the uniform lower form bound. Thus \(f_\e=\Phi_{b_\e}\circ G_\e\)
is balanced of degree one, converges strongly to \(F\), and
has \(\Yf\to0\) along this family.
\end{proof}

When the balanced zero supremum is attained, the maximizing map
is an orientation-preserving homothetic diffeomorphism for the
supplied metric. In fact \eqref{eq:sign} first makes it conformal.
Equality in \eqref{eq:balance-chain} then makes it a first eigenmap
with \(s=3\abs{df}^2\). On the volume-one round sphere one can instead
choose a reflection as \(F\) in the converse and insert two positive
bubbles. The maximizing sequence has degree one even though its
strong \(W^{1,2}\) limit has degree \(-1\).

\subsection{The critical energy required to lose degree}

Let \(\Omega\) be the oriented unit-round volume form on \(S^4\)
and \(v_4=\int_{S^4}\Omega\). The pointwise determinant inequality
\(\abs{J_f}\le\abs{df}^4/16\) is the arithmetic--geometric mean
inequality on the squared singular values. Thus a degree-one map
has fourth energy at least \(16v_4\). A stronger estimate applies
to the defect of a strongly converging sequence.

\begin{theorem}\label{thm:degree}
Let \(f_j:M\to S^4\) be smooth degree-one maps on a fixed
closed oriented Riemannian four-manifold, with bounded fourth
energy and strong ambient \(W^{1,2}\) limit a smooth map \(F\).
Along any subsequence for which the energy and signed Jacobian
measures converge weakly, write their limits as
\(\abs{dF}^4\,\dv+\nu\) and \(F^*\Omega+\eta\). Then
\begin{equation}\label{eq:atoms}
 \eta=v_4\sum_{\ell=1}^{N}k_\ell\delta_{p_\ell},\qquad
 k_\ell\in\mathbb Z\setminus\{0\},\qquad
 16\abs{\eta}\le\nu,\qquad
 \sum_\ell k_\ell=1-\deg F .
\end{equation}
In particular
\begin{equation}\label{eq:energy-gap}
 \liminf_j\int_M\abs{df_j}^4\,\dv
 \ge \int_M\abs{dF}^4\,\dv+16v_4\abs{1-\deg F}.
\end{equation}
If every atom of the limiting energy measure has mass strictly
below \(16v_4\), then \(\deg F=1\).
\end{theorem}

\begin{proof}
Put \(G_j=df_j-dF\). Interpolation gives \(G_j\to0\) in \(L^p\)
for \(p<4\). Expanding the fourth power shows
\(\abs{df_j}^4-\abs{dF}^4-\abs{G_j}^4\to0\) in \(L^1\),
so \(\nu\) is the nonnegative limit of \(\abs{G_j}^4\,\dv\).
The determinant expansion in an oriented orthonormal frame gives
\[
 J_{f_j}-J_F
 =\det(f_j,G_j(e_1),G_j(e_2),G_j(e_3),G_j(e_4))+o_{L^1}(1).
\]
All other terms have at most three \(G_j\) columns, or the bounded
columns of \(dF\) and the error \(f_j-F\).
Projecting the \(G_j\) columns onto \(f_j^\perp\) bounds the
displayed determinant by \(\abs{G_j}^4/16\). Taking limits
gives \(16\abs{\eta}\le\nu\) as measures.

For the integer statement fix a center \(x\) and geodesic spheres
of radius below the injectivity radius. On a subsequence with
summable squared \(L^2\) errors, coarea gives \(L^2\) convergence
on almost every sphere. Fatou gives a bounded tangential \(W^{1,4}\)
subsequence on each such sphere, since the bulk fourth energies are
bounded. Morrey compactness on the three-dimensional sphere gives
uniform convergence there. Exclude radii at which the limiting
energy has boundary mass. At a remaining radius the signed
Jacobians converge on the ball by the energy bound on boundary
neighborhoods. Join \(f_j|_{\partial B}\) to \(F|_{\partial B}\)
by the shortest geodesic homotopy. Its swept volume is bounded by
\[
 C\sup_{\partial B}\dist(f_j,F)
   \int_{\partial B}(\abs{d_\tau f_j}+\abs{d_\tau F})^3\,\dd A
 \longrightarrow 0 .
\]
The map on the double of the ball with this joining collar gives
an integer multiple of \(v_4\) by the degree formula. It can be
made smooth by making the ball maps constant in the normal
direction near the seams through a homotopy relative to the
boundary, which preserves their integrals. Hence
\(\eta(B_r(x))\in v_4\mathbb Z\) for good radii. Continuity from
below extends this to all small open balls. This is the local
degree argument underlying integer concentration for critical
Sobolev maps. See Hardt--Rivi\`ere~\cite{HardtRiviere2008Connecting}
for related concentration results.

Shrinking the balls gives integer point masses. Away from the
finitely many nonzero atoms, sufficiently small balls have total
variation \(<v_4\) and so signed mass zero. On every compact
set avoiding the atoms these radii can be chosen uniformly
by a finite covering. Averaging a continuous test over the small
balls then yields zero pairing with \(\eta\) on that set, by
Fubini and convergence of the normalized ball averages.
Thus the signed defect is supported on the stated atoms.
The degree integral gives the sum in \eqref{eq:atoms}.
The bound on total energy follows along a subsequence realizing
the liminf. The local cap excludes all the nonzero integers.
\end{proof}

A sufficient sequence-level cap is
\(\lim_{r\downarrow0}\limsup_j\int_{B_r(x)}\abs{df_j}^4\,\dv<16v_4\)
at every \(x\), since the energy atom of any extracted limit
obeys this bound. The coefficient is sharp. Each inserted conformal bubble
in the recovery of Theorem~\ref{thm:balance} contributes fourth
energy \(16v_4+o(1)\) by the integral of \(16b_\e^4\).
The cutoff and flattened portions contribute \(o(1)\), and the
small final conformal correction preserves this limit. The
remaining energy converges to the energy of \(F\). Thus the
recovery realizes the coefficient in \eqref{eq:energy-gap},
including the balanced reflection example. The energy defect
may additionally have a diffuse part, even though the signed
Jacobian defect is atomic.

\section{A sharp spectral criterion for a round Yamabe metric}\label{sec:rigidity}

We now drop the spin and orientation assumptions. Let \(g\) be any
smooth volume-one Yamabe minimizer on a closed connected Riemannian
four-manifold. When \(s=3\lambda_1\), the conformal second variation
vanishes in every first-eigenfunction direction. The fourth variation
at that degenerate minimum gives the following rigidity.

\begin{theorem}\label{thm:spectral}
With the positive eigenvalues counted with multiplicity,
\[
 A(g):=\frac35\sum_{i=1}^5\lambda_i(g)-s\ \ge 0 .
\]
Equality holds exactly when \((M,g)\) is homothetic to the round
four-sphere. At \(s=3\lambda_1\) the first eigenspace has dimension at
most five, and dimension five suffices for equality. Its orthonormal
basis, rescaled by \(1/\sqrt{5}\), gives a homothetic diffeomorphism to
the unit sphere.
\end{theorem}

\begin{proof}
The second variation in \eqref{eq:balance-chain} gives \(s\le3\lambda_1\),
so it suffices to study \(s=3\lambda\), \(\lambda=\lambda_1>0\).
For every real Sobolev test \(u\), the minimizing inequality is
\begin{equation}\label{eq:sharp-yamabe}
 \left[\int_M\left(u^2+\tfrac2\lambda\abs{\dd u}^2\right)\dv\right]^2
 \ge \int_M u^4\,\dv .
\end{equation}
The inequality extends to signed \(u\) by taking absolute values
in the Sobolev form. Let \(\phi\) be a real first eigenfunction with
\(\int\phi^2=1\). For \(w\) smooth of mean zero put
\(u=1+t\phi+t^2w\). The bracket in \eqref{eq:sharp-yamabe} is
\(1+3t^2+6t^3\int\phi w+t^4\int(w^2+2\abs{\dd w}^2/\lambda)\).
Writing \(z=t\phi+t^2w\), the right side is
\(1+6\int z^2+4\int z^3+\int z^4\).
Subtracting gives the cubic term \(-4t^3\int\phi^3\) and
the quartic term \(t^4\mathcal C_\phi(w)\), where
\begin{equation}\label{eq:quartic}
 \mathcal C_\phi(w)
 =9-\int\phi^4+\frac4\lambda\int\abs{\dd w}^2
            -4\int w^2-12\int\phi^2 w .
\end{equation}
The cubic term vanishes by the two signs of \(t\), and
\(\mathcal C_\phi(w)\ge0\) for every mean-zero \(w\).
Polarization shows that all triple products of first
eigenfunctions integrate to zero. The cubic degeneracy calculation
is related to that of Carlotto, Chodosh, and Rubinstein
\cite[Proposition~7]{CarlottoChodoshRubinstein2015Yamabe}.

Let \(b=\phi^2-1\) and \(q=\int\phi^4\).
Testing \(\Delta\phi=\lambda\phi\) with \(\phi^3\) gives
\(\int\abs{\dd b}^2=4\lambda q/3\), and
\(\int b^2=\int\phi^2 b=q-1\). Hence
\begin{equation}\label{eq:quartic-sharp}
 \mathcal C_\phi(b)=25-\tfrac{35}{3}q\ge0,\qquad
 q\le\tfrac{15}{7}.
\end{equation}
If equality holds, \(b\) minimizes the nonnegative functional
\(\mathcal C_\phi\), so its first variation against each mean-zero
test is zero. From \eqref{eq:quartic} this is
\(\Delta b=(5\lambda/2)b\). The product rule now gives
\[
 \abs{\dd\phi}^2=\tfrac{\lambda}{4}(5-\phi^2).
\]
Conversely the centered-square equation implies saturation by
testing with \(b\).

For an orthonormal first basis \(\phi_1,\ldots,\phi_d\), put
\(S=\sum_{i=1}^d\phi_i^2\). A unit coefficient vector \(a\),
averaged uniformly on \(S^{d-1}\), has
\(\mathbb E_a(a\cdot z)^4=3\abs z^4/(d(d+2))\). This follows
from rotation invariance and \(\sum_i a_i^2=1\).
Use \eqref{eq:quartic-sharp} for every direction \(a\) and
\(\int S=d\) to obtain
\[
 \frac{3d}{d+2}\le\frac{3}{d(d+2)}\int S^2\,\dv\le\frac{15}{7}.
\]
It follows that \(d\le5\). When \(d=5\) both inequalities
are equalities, so \(S=5\) everywhere and every direction has
quartic equality. Let \(X=(\phi_1,\ldots,\phi_5)/\sqrt5\).
By polarization of the gradient identity,
\begin{equation}\label{eq:point-gram}
 \ip{\dd X_i}{\dd X_j}_g
       =\tfrac{\lambda}{4}(\delta_{ij}-X_iX_j),\qquad
 X^*h_{\rd}=\tfrac{\lambda}{4}g .
\end{equation}
Indeed \(X\) is sphere-valued by \(S=5\), and the matrix on
the right in the first formula is \(\lambda/4\) times the
orthogonal tangent projection. The four nonzero singular
values of \(dX\) are therefore equal to \(\sqrt{\lambda/4}\).
The resulting local homothety is a covering by compactness,
and the target is simply connected. This derives the metric
from the quartic Yamabe inequality through the five coordinates.
Obata's classical sphere theorem uses a prescribed
Hessian equation for one function~\cite[Theorem~A]{Obata1962Sphere}.

For the round converse, if the sectional curvature is \(K\),
then \(s=12K\) and the five coordinate eigenfunctions have
eigenvalue \(4K\). The integrated Bochner identity with
\(\Ric=3Kg\) and \(\abs{\nabla^2 u}^2\ge(\Delta u)^2/4\)
gives \(\lambda_1\ge4K\). Hence the first five eigenvalues
have the stated value and \(A(g)=0\).
\end{proof}

If \(F:M\to S^4\) is a first eigenmap at \(s=3\lambda_1\),
its image spans \(\RR^5\) exactly when its five coordinates
are independent. In that case the dimension is five and
\(F=B X\) for some real \(5\times5\) matrix \(B\).
Since \(X\) is surjective and \(\abs{BX}=1\), we have
\(B^{\mathsf T}B=I\). Thus every full-span first eigenmap in
this setting is a homothetic diffeomorphism. In particular a
degree-zero limit in Theorem~\ref{thm:balance} must lie in a
proper great subsphere.

\subsection{Simultaneous normalization of mean and covariance}
\label{sec:moment-normalization}

For a smooth degree-one map \(f\) from a closed connected oriented four-manifold
with a smooth positive probability volume form \(\dd v\),
we can impose
\[
 \int f\circ T\,\dd v=0,\qquad
 \int (f\circ T)(f\circ T)^{\mathsf T}\,\dd v=I/5
\]
by an orientation-preserving diffeomorphism \(T\) of the domain.
To prove this, consider
\(\Phi(y)=(y,yy^{\mathsf T}-I/5)\) in
\(\RR^5\oplus\operatorname{Sym}_0(5)\). The image has full affine
span because a trace-free quadratic plus a linear function
constant on the unit sphere must be zero.
Its round mean is zero. Thus every nonzero linear functional
on this space takes both signs on the image. The measure
\(\mu=f_*\dd v\) has full support by nonzero degree and positivity
of the volume form. Define
\[
 \mathcal L(b)=\log\int_{S^4}\exp\ip{b}{\Phi(y)}\,\dd\mu(y).
\]
Its derivative is the tilted mean and its Hessian is the tilted
covariance, positive definite by full support and the affine span.
For every unit \(b\), the maximum of \(\ip{b}{\Phi(y)}\) is positive.
Compactness of the unit parameter sphere gives a uniform positive
lower bound on these maxima. Finite neighborhoods in parameter
space and full support then give
\(\mathcal L(tb)\ge c t+\log m\) for constants \(c,m>0\) and \(t\ge0\).
Consequently a unique minimizing tilt gives mean zero for \(\Phi\).
It has the smooth positive density
\(\rho(x)=\exp(\ip{b_*}{\Phi(f(x))}-\mathcal L(b_*))\)
on the domain. Moser transport gives \(T^*(\rho\,\dd v)=\dd v\)
\cite{Moser1965Volume}. Indeed, on a closed connected oriented
manifold with equal total volume, the difference of the volume
forms is an exact top form \(\dd\alpha\). Along their linear
interpolation \(\omega_t\), the flow of the time-dependent field
satisfying \(\iota_{V_t}\omega_t=-\alpha\) has constant pulled back
volume form. At \(t=1\) this is the claimed transport. It preserves
degree and gives the two stated moments. The construction
establishes feasibility of the constraints for each map, with
derivative bounds depending on the map and the transport.

\subsection{A covariance bound for the coupled deficit}

For a balanced map \(f\) on the supplied Yamabe minimizer, set
\(C_f=\int ff^{\mathsf T}\,\dv\), so \(\tr C_f=1\).
Suppose \(C_f\ge\delta I\) for \(0<\delta\le1/5\).
The coordinates of \(\varphi=C_f^{-1/2}f\) have mean zero
and are orthonormal in \(L^2\). The matrix
\(G_{ij}=\int\ip{\dd\varphi_i}{\dd\varphi_j}\,\dv\)
satisfies \(G\ge\lambda_1 I\) by Poincar\'e,
and \(\tr G\ge\sum_{i=1}^5\lambda_i\) by the variational trace
principle. For the latter, expand the orthonormal functions in a
Laplace eigenbasis. The summed squared coefficient at each mode
is in \([0,1]\), and all weights sum to five.
Consequently
\[
 \int\abs{df}^2\,\dv=\tr(C_fG)
 \ge \lambda_1+\delta\sum_{i=2}^5(\lambda_i-\lambda_1).
\]
Combining with the scalar constant test gives
\begin{equation}\label{eq:covariance-gap}
 \Yf\le s-3\lambda_1-3\delta\sum_{i=2}^5(\lambda_i-\lambda_1).
\end{equation}
On a fixed nonround minimizing metric, this bound is strictly
negative by Theorem~\ref{thm:spectral}. In particular the supremum
over balanced degree-one maps with \(C_f=I/5\) is zero exactly for
a homothetically round metric, where an oriented homothety attains
it. If \(f_j\) is balanced with \(Y_{f_j}([g])\to0\) on a fixed nonround
minimizer, its least covariance eigenvalue must tend to zero. The
full covariance criterion is independent of the Jacobian criterion
in Theorem~\ref{thm:degree}.

\subsection{Local minimizers and conformal vector fields}\label{sec:local}

The proof of Theorem~\ref{thm:spectral} uses the minimizing inequality
only along curves \(u_t=1+t\phi+t^2w\) with \(t\to0\). Call a smooth
volume-one metric \(g\) of constant scalar curvature \(s\) a
\emph{local Yamabe minimizer} if, for all smooth \(\phi\) and \(w\),
the Yamabe quotient
\(\int_M(6\abs{\dd u_t}^2+su_t^2)\,\dv/\norm{u_t}_{4,g}^2\)
is at least \(s\) for all sufficiently small \(\abs t\).
Yamabe minimizers have this property, and so do local minima of the
Yamabe quotient of the class in the smooth topology.

\begin{proposition}\label{prop:local}
Let \(g\) be a local Yamabe minimizer on a closed connected
four-manifold. Then \(s\le3\lambda_1\) and \(A(g)\ge0\), and every
statement of Theorem~\ref{thm:spectral} holds for \(g\).
\end{proposition}

\begin{proof}
The expansion of the quotient of \(1+t\phi\) used for
\eqref{eq:balance-chain} gives \(s\le3\lambda_1\). Hence
\(A(g)\ge3\lambda_1-s\ge0\), and \(A(g)>0\) unless \(s=3\lambda_1>0\).
At \(s=3\lambda_1\) the inequality \eqref{eq:sharp-yamabe} holds for
\(u=u_t\) and small \(\abs t\), for each first eigenfunction \(\phi\)
and each smooth \(w\) of mean zero. The proof of
Theorem~\ref{thm:spectral} uses it only for these tests. The cubic
and quartic coefficients give the vanishing triple products and
\(\mathcal C_\phi(w)\ge0\). The equality case uses the first variation
of \(\mathcal C_\phi\) at the smooth function \(\phi^2-1\), and the
remaining steps are pointwise.
\end{proof}

Let \(X\) be a conformal vector field for a metric \(g\) of constant
scalar curvature \(s\) on a closed four-manifold, and write
\(\mathcal L_Xg=2\sigma g\), so that
\(\sigma=\tfrac14\operatorname{div}X\). Its flow \(\Phi_\tau\) satisfies
\(\Phi_\tau^*g=a_\tau^2g\) with \(a_\tau=1+\tau\sigma+O(\tau^2)\), and
\(\Phi_\tau^*g\) has scalar curvature \(s\). By
\eqref{eq:conformal-L} without the map potential,
\(s\,a_\tau^3=6\Delta_ga_\tau+s\,a_\tau\). The first-order term in
\(\tau\) is
\begin{equation}\label{eq:conformal-eigen}
 \Delta_g\sigma=\tfrac{s}{3}\sigma .
\end{equation}

\begin{proposition}\label{prop:broken}
Let \(g\) be a local Yamabe minimizer on a closed connected
four-manifold, not homothetic to the round sphere. If \(g\) has a
conformal vector field that is not Killing, then \(s=3\lambda_1>0\).
The quotient of the Lie algebra of conformal vector fields of \(g\) by
its Killing fields has dimension at most four.
\end{proposition}

\begin{proof}
A conformal field with \(\sigma=0\) is Killing, so \(X\mapsto\sigma\)
is injective on the quotient. Its image lies in the eigenspace of
\(s/3\) by \eqref{eq:conformal-eigen}. A nonzero \(\sigma\) has mean
zero and is not constant, so \(s/3\) is a positive eigenvalue, and
Proposition~\ref{prop:local} gives \(s/3=\lambda_1\). The first
eigenspace has dimension at most five, and dimension five would make
\(g\) round.
\end{proof}

The Delaunay metrics of Theorem~\ref{thm:necklace} are minimizers with
a conformal field that is not Killing, so the scalar equality
\(s=3\lambda_1\) does occur on nonround minimizers.

\section{Lower-dimensional eigenmaps and metric variations}
\label{sec:lower}

The degree-zero alternatives at \(s=3\lambda_1\) lie in lower
great spheres by Theorem~\ref{thm:spectral}. We record the
geometric restrictions supplied by the fixed-metric equations.
In this section \(M\) is a smooth oriented homotopy four-sphere,
\(g\) is a volume-one Yamabe minimizer at \(s=3\lambda_1(g)\),
\(\lambda=\lambda_1(g)\), and \(F\) is a supplied first eigenmap.

\subsection{Rank and the circle Hopf bundle}

The equation \(\Delta_g F=\lambda F\), \(\abs F=1\) gives
\(\abs{dF}^2=\lambda>0\). On an open set of rank one the
constant-rank theorem gives \(F=\gamma\circ u\) for a local
unit-speed curve \(\gamma\). Then \(\abs{\dd u}^2=\lambda\)
and
\(\gamma'\Delta_g u-\gamma''\abs{\dd u}^2=\lambda\gamma\).
The tangent part gives \(\Delta_g u=0\), so \(\gamma''=-\gamma\)
and the image there lies in a great circle. Unique continuation
for the scalar eigenfunctions implies that the image lies in
that circle globally~\cite{Aronszajn1957Continuation}.
A global angle lift exists by simple connectivity, and its
harmonic equation on the closed manifold makes it constant.
Thus rank at least two occurs on an open dense set. A map
to a great \(S^2\) cannot be a submersion everywhere. Its
kernel and orthogonal complement would be oriented two-plane
bundles with zero Euler classes because \(H^2(M;\mathbb Z)=0\),
contradicting \(\chi(M)=2\). For a great \(S^3\), an everywhere
submersion would give an oriented tangent line subbundle and
a nowhere-zero vector field, with the same contradiction.

For a map to a great \(S^2\), pull back the circle Hopf bundle and
choose a smooth lift \(z:M\to S^3\subset\mathbb C^2\), possible
because the first Chern obstruction is zero. With the unit-radius
target, the Hopf connection \(a=\operatorname{Im}(z^*\dd z)\)
satisfies
\begin{equation}\label{eq:circle-hopf}
 \abs{\dd z}^2=\lambda/4+\abs a^2,\qquad
 \dd a=\tfrac12F^*\Omega_2 ,
\end{equation}
where \(\Omega_2\) is the unit sphere area form with the standard
Hopf orientation. At \(z=(1,0)\) the differential on the complex
horizontal line doubles lengths and the exterior derivative of
the connection has value \(2\) on its oriented orthonormal basis.
The pulled-back area has value \(4\). Equivariance gives
\eqref{eq:circle-hopf} everywhere. Change the phase by solving
\(\Delta_g\varphi=-d^*a\) so that \(a+\dd\varphi\) is coclosed.
It is then coexact because \(H^1(M;\RR)=0\). Use this lift below.

Set \(m=\int z\,\dv\). The Hopf formula expresses \(zz^*\)
as \(I/2\) plus a trace-free matrix linear in \(F\). Balance gives
\(\int zz^*\,\dv=I/2\), and
\(\int\abs{v^*z-v^*m}^2\,\dv=1/2-\abs{v^*m}^2\) for a unit
\(v\in\mathbb C^2\). Taking \(v=m/\abs m\) when \(m\ne0\) gives
\(\abs m^2\le1/2\). At equality, \(v^*z\) would be constant with
modulus \(1/\sqrt2\), so the Hopf image would be orthogonal to
the Hopf image of \(v\). The resulting great-circle map was
excluded above. Hence \(\abs m^2<1/2\).
Scalar Poincar\'e and \eqref{eq:circle-hopf} now give
\[
 \norm a_2^2\ge\lambda(3/4-\abs m^2)>\lambda/4,\qquad
 \norm{\dd a}_2^2
 \le \lambda^2/16 ,
\]
since the product of the two singular values of \(dF\) is at
most \(\lambda/2\). The Rayleigh quotient of \(a\) bounds the
first positive coexact one-form eigenvalue by
\[
 \lambda_{\mathrm{coex}}^{(1)}
 \le \frac{\lambda}{16(3/4-\abs m^2)}
 <\lambda/4=s/12 .
\]
The last equality uses \(s=3\lambda\). The area pullback is a
nonzero exact decomposable form, so its two chiral parts have
equal norm pointwise, and it cannot be coclosed by integration
by parts. The obstruction here requires an independent coexact
gap to exclude the target on a given metric. For comparison, the
round product \(S^2(a)\times S^2(a/\sqrt2)\) has \(s=3\lambda_1\)
and a rank-two first eigenmap by projection to the first factor.
Its pulled-back Hopf class is nonzero on that factor, as is
seen by integrating \(F^*\Omega_2\). This comparison uses
a supplied constant-scalar product metric
\cite[\S4.2]{DeLimaPiccioneZedda2012Yamabe}.

\subsection{The quaternionic lift of a degree-zero map}

The Hopf lift for \(S^4\cong\mathbb{HP}^1\) gives a related
bound for the great \(S^3\) alternative. For a unit quaternionic
column \(z\) represent the target point by
\[
 P=zz^*=\tfrac12(I+T(F)),\qquad
 T(q,t)=
 \begin{pmatrix}t&q\\\bar q&-t\end{pmatrix},\qquad
 q=2z_1\bar z_2,\quad t=\abs{z_1}^2-\abs{z_2}^2 .
\]
The tautological quaternionic line has top Chern number of
absolute value one, by the Hopf clutching map of degree one.
On a four-dimensional base the pulled-back unit sphere bundle
has a section exactly when its top Chern class vanishes. The
fiber \(S^3\) is two-connected, so its only obstruction is in
degree four~\cite[\S3]{Hatcher2017VectorBundles}. Thus
\(\deg F=0\) is precisely the condition for a lift.

For any smooth lift put \(a=z^*\dd z\), \(b=(I-P)\dd z\).
The relation \(T(v)^2=\abs v^2 I\) gives
\(\abs{b(X)}=\abs{dF(X)}/2\), hence
\[
 \abs{\dd z}^2=\lambda/4+\abs a^2,\qquad
 F_a=\dd a+a\wedge a=b^*\wedge b.
\]
In a unit frame of the orthogonal quaternionic line the last
form acts as \(2\operatorname{Im}(\bar b(X)b(Y))\), with norm
\(\abs{dF(X)\wedge dF(Y)}/2\).
In particular \(\abs{F_a}^2\le 3\lambda^2/32\) by the same
singular-value inequality used in Section~\ref{sec:twisted}.
Equality everywhere would force a full-rank map, incompatible
with degree zero. As in the complex case,
\(\int P\,\dv=I/2\) and \(\abs{\int z\,\dv}^2\le1/2\).
Poincar\'e gives \(\int\abs a^2\,\dv\ge\lambda/4\).
If \(F\) has full linear span this lower bound is strict
uniformly over lifts of this fixed map. In fact
\[
 \abs{\int z\,\dv}\le \max_{\abs v=1}\int\abs{Pv}\,\dv
                 <1/\sqrt2.
\]
Equality in the last comparison would require \(\abs{Pv}\) to
be constant, which by the projector formula confines \(F\) to a
great \(S^3\). Conversely if the image lies there and \(v\) is a
unit column projecting to its normal, \(z=\sqrt2 Pv\) is a unit
lift with mean \(v/\sqrt2\). It is affine-linear in \(F\), so it
attains the Poincar\'e inequality and \(\int\abs a^2\,\dv=\lambda/4\).
The connection equation now includes \(a\wedge a\), so the
linear coexact argument for \eqref{eq:circle-hopf} leaves a
separate gauge estimate to be established.

\subsection{Bochner and stress identities at scalar equality}

At \(s=3\lambda\), let \(B=F^*h_{\rd}\), \(H=B-\lambda g/4\),
and let \(D_0=\abs{\nabla dF}^2\) use the target Levi--Civita
connection. The eigenmap is harmonic with constant energy.
Commuting the two covariant derivatives of \(dF\) gives the
Bochner identity
\begin{equation}\label{eq:map-bochner}
 0=-\tfrac12\Delta_g\abs{dF}^2
   =D_0+\ip{\Ric}{B}+\abs B^2-\lambda^2 ,
 \qquad D_0+\abs H^2=-\ip{\Ric_0}{H}.
\end{equation}
The target-curvature contribution is
\(\sum_{i,j}(B_{ij}^2-B_{ii}B_{jj})\) by round curvature one.
The trace terms in the Ricci and pullback contractions cancel
\(\lambda^2\).
At rank \(k\), \(\abs B^2\ge\lambda^2/k\), so
\[
 \Ric_{\min}\le (1-1/k)\lambda-D_0/\lambda,\qquad
 \abs{\Ric_0}\ge\lambda\sqrt{1/k-1/4}\quad(k<4).
\]
At rank one with \(B=\lambda u^\flat\otimes u^\flat\), its
Ricci direction satisfies \(\Ric(u,u)=-D_0/\lambda\).
Positive Ricci therefore excludes the compulsory rank-one
point for a great \(S^2\). A lower bound strictly above
\(\lambda/2\) excludes the compulsory rank drop for a great
\(S^3\). Full rank everywhere gives a covering and hence
a diffeomorphism.

There is also a finite stress criterion, with the map held fixed
when varying the metric. Put \(T_0=\Ric-3B=\Ric_0-3H\).
The tensors \(H,T_0\) are trace-free and divergence-free.
The divergence of \(B\) is the tension contraction plus
\(\tfrac12\dd\abs{dF}^2\), and the scalar curvature is constant.
The ground state of \(L_{g,F}\) at this pair is the constant
simple eigenfunction with eigenvalue zero, because \(s=3\lambda\).
Along \(g+tK\), differentiating the smooth simple eigenpair
and integrating at \(t=0\) gives
\begin{equation}\label{eq:stress}
 \dot\ell_F[K]=-\int_M\ip{T_0}{K}\,\dv,\qquad
 \dot\ell_F[H]=\int_M(D_0+4\abs H^2)\,\dv .
\end{equation}
Indeed, the metric derivative of the Laplacian annihilates
constants. The derivative of the scalar term is
\(-\ip{\Ric}{K}+\nabla^j\nabla^i K_{ij}+
\Delta_g\tr_g K\), whose last terms integrate to zero.
The derivative of \(-3\abs{dF}^2\) is \(3\ip{B}{K}\).
Equation~\eqref{eq:map-bochner} then gives the second identity.
The test \(H\) has zero first volume variation. Its derivative
vanishes only for \(B=(\lambda/4)g\), yielding a homothetic
diffeomorphism.

For a sequence \(f_j\to F\) strongly in \(W^{1,2}\) on a fixed
\(g\), the same criterion can be tested by finitely many moments.
Choose a real orthonormal basis \((\phi_a)\) of the first
eigenspace and take
\(K_{ab}=(\dd\phi_a\odot\dd\phi_b)_0\), where
\(\odot\) denotes the symmetric product and the subscript is
the trace-free part. If every pairing of
\(\Ric_0-3(f_j^*h_{\rd})_0\) with \(K_{ab}\) tends to zero,
the pullback tensors converge in \(L^1\) by strong derivative
convergence. Thus the pairings of \(T_0\) with all \(K_{ab}\)
vanish. Since \(H\) is a linear combination of the \(K_{ab}\)
by the first eigenfunction expansion of \(F\),
\eqref{eq:stress} gives homothety. Obtaining those metric
stationarity conditions from a geometrically constructed
sequence remains a separate premise.

\Needspace{15\baselineskip}
\section{Relative pinching and the joint law of five modes}
\label{sec:near}

Consider smooth volume-one Yamabe minimizers, now allowing the
manifold and metric to vary. Assume \(s>0\), and write
\[
 \e=\frac{3\lambda_5}{s}-1,\qquad
 \frac{\e}{5}\le \frac{A(g)}{s}\le\e
\]
by \(s\le3\lambda_1\) and the monotonicity of the eigenvalues.
The scale in this comparison is essential when \(s\) can tend to
zero.

\begin{theorem}\label{thm:near}
Let \(u_1,\ldots,u_5\) be an \(L^2\)-orthonormal family of the first
five positive Laplace eigenfunctions, in eigenvalue order,
and put \(X=(u_1,\ldots,u_5)/\sqrt5\).
There is a universal bound on \(\abs X\) for \(0\le\e\le1\).
For a sequence of such minimizers with \(\e\to0\),
the probability measures \(X_*(\dd V_g)\) on \(\RR^5\) converge
weakly to uniform probability on the unit \(S^4\).
In the scale
\[
 L=\tfrac3s\Delta_g,\qquad
 \Gamma(a,b)=\tfrac3s\ip{\dd a}{\dd b}_g,\qquad
 \norm{w}_{\mathcal H}^2=\int_M(w^2+\Gamma(w,w))\,\dv ,
\]
one has
\begin{align}
 \norm{\abs X^2-1}_{2,g}&\le C\e^{1/4},\notag\\
 \left\|\Gamma(X_i,X_j)-\tfrac14(\delta_{ij}-X_iX_j)\right\|_{\mathcal H^*}
 &\le C\e^{1/4}.
 \label{eq:near-gram}
\end{align}
\end{theorem}

\begin{proof}
It suffices to treat \(0<\e\le1\), since the case \(\e=0\)
is Theorem~\ref{thm:spectral}. Define
\(\mathcal Q(v)=\int(v^2+2\Gamma(v,v))\,\dv\).
By the minimizing inequality,
\begin{equation}\label{eq:near-sob}
 \norm v_4^2\le\mathcal Q(v)\le2\norm v_{\mathcal H}^2.
\end{equation}
This holds for signed \(v\) by taking absolute values.
The eigenvalues \(\nu_i\) of \(L\) on the \(u_i\) satisfy
\(1\le\nu_i\le1+\e\). For a normalized eigenfunction \(u_i\),
testing with \(u_i^{p-1}\) for dyadic \(p\ge2\) gives
\(\int\Gamma(u_i^{p/2},u_i^{p/2})=
\nu_i p^2\norm{u_i}_p^p/(4(p-1))\).
Applying \eqref{eq:near-sob} to \(u_i^{p/2}\) yields
\[
 \norm{u_i}_{2p}\le
 \left(1+\frac{\nu_i p^2}{2(p-1)}\right)^{1/p}\norm{u_i}_p .
\]
The product over dyadic \(p\) converges uniformly in
\(\nu_i\le2\). Thus all five modes have a universal
\(L^\infty\) bound.

For a unit coefficient vector take \(\phi=\sum_i a_i u_i\).
Then \(\int\phi=0\), \(\norm{\phi}_2=1\),
\(r=(L-1)\phi\) has \(\norm r_\infty\le C\e\), and
\(\alpha=\int\phi r\in[0,\e]\).
For a smooth mean-zero \(w\) use the two tests
\(v_\pm=1\pm t\phi+t^2w\).
Write \(q=\int\phi^4\), \(a_0=3+2\alpha\),
\(J_w=\mathcal Q(w)\) and \(B_w=\int(\phi w+2\Gamma(\phi,w))\).
The sign average of the nonnegative deficits
\(\mathcal Q(v_\pm)^2-\int v_\pm^4\) is
\begin{align}
 4\alpha t^2
 &+\bigl(\mathcal C_\phi(w)+12\alpha+4\alpha^2\bigr)t^4\notag\\
 &+\left(2a_0J_w+4B_w^2-4\int w^3-6\int\phi^2w^2\right)t^6
  +\left(J_w^2-\int w^4\right)t^8 ,
 \label{eq:averaged-quartic}
\end{align}
where \(\mathcal C_\phi(w)=9-q+4\int\Gamma(w,w)
-4\int w^2-12\int\phi^2w\).
Indeed, the quadratic bracket for the two signs is
\(1+a_0t^2\pm2B_wt^3+J_wt^4\), and the average of the
fourth powers uses \((1+t^2w)^4+6t^2\phi^2(1+t^2w)^2+t^4\phi^4\).
The last two coefficients in \eqref{eq:averaged-quartic} are
bounded for \(\norm w_{\mathcal H}\le K\), by
\eqref{eq:near-sob}, Cauchy--Schwarz, and the sup norm of \(\phi\).
With \(t=\e^{1/4}\) this proves
\(\mathcal C_\phi(w)\ge-C_K\sqrt\e\) uniformly on these bounded sets.

Put \(b=\phi^2-1\). The chain rule gives
\(\int\Gamma(b,b)=\frac43\int\phi^3L\phi\).
It follows that \(\norm b_{\mathcal H}\) is uniformly bounded and
\[
 \mathcal C_\phi(b)=25-\tfrac{35}{3}q+\tfrac{16}{3}\int\phi^3r .
\]
In particular \(q\le15/7+C\sqrt\e\).
Average over the coefficient sphere as in the exact proof, now
with \(S=\sum_i u_i^2\). The average is
\((3/35)\int S^2=15/7+(3/35)\int(S-5)^2\).
This implies the radial estimate in \eqref{eq:near-gram}.
It also gives a lower bound in every direction. Indeed,
\((15/7+C\sqrt\e)\abs a^4-q(a)\) is a nonnegative homogeneous
quartic on the unit coefficient sphere and has
\(L^1\) norm \(O(\sqrt\e)\). Equivalence of the \(L^\infty\) and
\(L^1\) norms on the finite-dimensional space of spherical quartics
then gives
\begin{equation}\label{eq:near-moment}
 \abs{q-15/7}\le C\sqrt\e
\end{equation}
uniformly in the direction \(a\).

Now let \(\zeta\) be smooth with mean zero and
\(\norm{\zeta}_{\mathcal H}\le1\). The change from \(b\) to
\(b+\tau\zeta\) in \(\mathcal C_\phi\) is
\[
 8\tau\int ((L-\tfrac52)b)\zeta\,\dv+
    4\tau^2\int(\Gamma(\zeta,\zeta)-\zeta^2)\,\dv .
\]
The value at \(b\) is \(O(\sqrt\e)\) by \eqref{eq:near-moment}.
Use the lower bound on \(\mathcal C_\phi(b+\tau\zeta)\) for the
two signs \(\tau=\pm\e^{1/4}\).
The quadratic term is \(O(\sqrt\e)\), giving
\(\norm{(L-\tfrac52)b}_{\mathcal H^*}\le C\e^{1/4}\).
This residual annihilates constants, so the estimate holds also
for tests with nonzero mean, by subtracting the mean.
The nonnegative Laplacian product rule gives
\[
 \Gamma(\phi,\phi)-\tfrac14(5-\phi^2)
       =\phi r-\tfrac12(L-\tfrac52)b .
\]
The \(L^2\) norm of \(\phi r\) is \(O(\e)\).
Polarization and division by five now give the dual Gram
estimate in \eqref{eq:near-gram}.

It remains to identify the law. The vectors \(X\) take values
in a common compact ball \(K\). For any fixed smooth \(F\)
of the target variables,
\[
 \norm{F(X)}_{\mathcal H}^2\le
  \sup_K\abs F^2+(1+\e)\sup_K\abs{\nabla_{\RR^5}F}^2 .
\]
In fact the pointwise Gram matrix \((\Gamma(X_i,X_j))\) is
positive semidefinite with integrated trace
\(\sum_i\nu_i/5\le1+\e\). Thus each \((\partial_jF)(X)\)
is a valid uniform test in the dual estimate.
By the chain rule and integration by parts for
\(LX_i=\nu_iX_i\), every limiting law \(\mu\) satisfies
\begin{equation}\label{eq:sphere-law}
 4\int x_i F(x)\,\dd\mu=
    \sum_j\int(\delta_{ij}-x_ix_j)\partial_jF(x)\,\dd\mu .
\end{equation}
The radial bound gives support on \(S^4\).
The tangent field \(e_i-x_i x\) on the unit sphere has
divergence \(-4x_i\), so round probability satisfies
\eqref{eq:sphere-law}. To check uniqueness, take
\(F=\partial_i P\) and sum in \(i\) for each homogeneous
polynomial \(P\) of degree \(m\ge1\).
Euler's identity gives
\[
 m(m+3)\int P\,\dd\mu=\int\sum_i\partial_{ii}P\,\dd\mu.
\]
The right side is a moment of degree \(m-2\), or zero.
Starting from total mass one, this determines all moments.
Polynomials are dense on the compact sphere, giving a unique
limit and hence convergence of the entire sequence.
\end{proof}

The estimates use only the scaled Yamabe inequality.
They bound the volume of the set where \(\abs X^2\) differs
from one by a fixed amount, and test the differential Gram
matrix in a dual energy norm. They leave pointwise rank and
the geometry of small-volume subsets uncontrolled.
Theorem~\ref{thm:necklace} gives volume-one minimizers on
\(S^1\times S^3\) with \(\e\to0\), so relative pinching alone does
not identify the manifold.

An additional Ricci lower bound gives a separate smooth
recognition criterion. Suppose
\(\Ric_g\ge(1-\delta)(s/4)g\) with \(0\le\delta<1\).
For \(k=(1-\delta)sg/12\), one has \(\Ric_k\ge3k\), and
\[
 \lambda_4(k)-4\le\frac{4(\e+\delta)}{1-\delta}.
\]
Aubry's sphere theorem assumes \(\Ric_k\ge3k\) in dimension four
and sufficiently small pinching of the fourth positive eigenvalue.
It gives diffeomorphism to the standard sphere and a
Gromov--Hausdorff estimate~\cite[Theorem~1]{Aubry2005Pincement}.
In particular this route applies for \(\e,\delta\to0\).
The bound on \(\Ric_g\) is an independent geometric premise.
Proposition~\ref{prop:budget} gives a second criterion on a homotopy
four-sphere, in which an integral bound on the trace-free Ricci
tensor replaces both the pointwise bound and the spectral pinching.

\section{A flat neck, Delaunay minimizers, and an integral curvature criterion}\label{sec:necklace-tests}
\label{sec:tests}

The construction problem asks for conformal classes on a homotopy
four-sphere whose Yamabe minimizers are relatively pinched, together
with enough geometric control to recognize the manifold. This section
records three exact statements that delimit it. Attaching a round
region across a conformally flat neck does not change the conformal
class. On \(S^1\times S^3\) the minimizers are explicit and become
relatively pinched as the circle lengthens. At a Yamabe metric the
hypothesis of Chang, Gursky, and Yang is a bound on the trace-free
Ricci tensor, and the minimizers on \(S^1\times S^3\) give equality in
it. Throughout, \(Y(\Sphere)=8\sqrt6\,\pi\) is the Yamabe constant of
the round class, so \(Y(\Sphere)^2=384\pi^2\).

\subsection{A round region across a conformally flat neck}
\label{sec:flat-neck}

Let \((M,g_0)\) be a closed Riemannian four-manifold that is flat on a
coordinate ball \(\{\abs y<2\rho\}\) centered at \(p\).
Proposition~\ref{prop:green} supplies such metrics on every smooth
homotopy four-sphere. Let \(z\) be a stereographic coordinate on the
unit round sphere, so that
\(h_{\rd}=4\abs{\dd z}^2/(1+\abs z^2)^2\), and let \(q^*\) be the point
\(z=\infty\). Fix \(\kappa>0\) and a reflection \(S\) of \(\RR^4\), and
put
\[
 \iota(y)=\frac{\kappa\,Sy}{\abs y^2},\qquad
 M_{\rho,\kappa}=\bigl(M\setminus\{\abs y\le\rho\}\bigr)\cup_{\iota}
        \bigl(\Sphere\setminus\{\abs z\le\kappa/(2\rho)\}\bigr),
\]
where \(\iota\) identifies the collar \(\{\rho<\abs y<2\rho\}\) with the
collar \(\{\kappa/(2\rho)<\abs z<\kappa/\rho\}\). The map \(\iota\) is
an inversion composed with a reflection and a dilation. It preserves
orientation and is conformal from \(\abs{\dd y}^2\) to
\(\abs{\dd z}^2\). Since \(g_0=\abs{\dd y}^2\) on the first collar and
\(h_{\rd}\) is a positive multiple of \(\abs{\dd z}^2\) on the second,
the two conformal classes agree under the identification. They define
a conformal class on \(M_{\rho,\kappa}\), which we write
\([g_{\rho,\kappa}]\). A partition of unity applied to the conformal
factors gives representatives that are homothetic to \(g_0\) on the
core \(M\setminus\{\abs y<2\rho\}\), round on
\(\{\abs z\ge\kappa/\rho\}\), and conformally flat on the neck between
them. The parameter \(\kappa\) fixes the size of the hole in the round
region, and the homothety factor on the core fixes the relative scale
of the two pieces.

\begin{proposition}\label{prop:flat-neck}
Define \(\Psi:M\to M_{\rho,\kappa}\) to be the identity on
\(M\setminus\{\abs y\le\rho\}\), the map \(y\mapsto\iota(y)\) on
\(\{0<\abs y<2\rho\}\), and \(\Psi(p)=q^*\). Then \(\Psi\) is a
diffeomorphism and \(\Psi^*[g_{\rho,\kappa}]=[g_0]\). Consequently
\(Y([g_{\rho,\kappa}])=Y([g_0])\) for all \(\rho\) and \(\kappa\), and
\(\Psi\) carries the volume-one Yamabe minimizers of \([g_0]\)
isometrically onto those of \([g_{\rho,\kappa}]\). Their scalar
curvature, their Laplace spectrum, and the quantities \(A\) and \(\e\)
do not depend on \(\rho\), on \(\kappa\), or on the representative
chosen on the neck.
\end{proposition}

\begin{proof}
On the collar the two definitions give points that the gluing
identifies, so \(\Psi\) is well defined. The formula for \(\iota\) maps
\(\{0<\abs y<2\rho\}\) diffeomorphically onto
\(\{\abs z>\kappa/(2\rho)\}\). In the chart \(w=z/\abs z^2\) at \(q^*\)
it reads \(w=Sy/\kappa\), which is linear, so the extension over \(p\)
is smooth with smooth inverse. The two pieces of \(M_{\rho,\kappa}\)
are the images of \(M\setminus\{\abs y\le\rho\}\) and of
\(\{\abs y<2\rho\}\), which gives bijectivity. On the first piece
\(\Psi\) is the identity. On the second it is conformal from
\(\abs{\dd y}^2\) to \(\abs{\dd z}^2\), hence to \(h_{\rd}\). The Yamabe
constant and the set of minimizing metrics are determined by the
conformal class, which gives the remaining statements.
\end{proof}

The minimizing metric therefore does not register the round region.
Let \(g=u^2g_0\) be a volume-one Yamabe minimizer of \([g_0]\), with
\(0<u_-\le u\le u_+\) on \(\{\abs y\le2\rho\}\). Under \(\Psi\) the
round region and the neck correspond to the ball
\(\{\abs y<2\rho\}\), whose \(g\)-volume is at most
\(8\pi^2u_+^4\rho^4\). Let \(\phi\) be a Lipschitz function on \(M\)
that is constant outside this ball, with mean \(\bar\phi\) for
\(\dv\). Then
\begin{equation}\label{eq:chamber-rayleigh}
 \int_M\abs{\dd\phi}_g^2\,\dv\ \ge\
 \frac{u_-^2}{u_+^4}\,\frac{\mu_1}{4\rho^2}
 \int_M(\phi-\bar\phi)^2\,\dv ,
\end{equation}
where \(\mu_1\) is the first Dirichlet eigenvalue of the unit ball of
\(\RR^4\). Indeed, let \(c\) be the constant value. The mean minimizes
the \(L^2\) distance to constants, so the integral on the right is at
most \(u_+^4\int_{\abs y<2\rho}(\phi-c)^2\,\dd y\). The function
\(\phi-c\) vanishes on the boundary of the ball, and
\(\abs{\dd\phi}_g^2\,\dv=u^2\abs{\dd\phi}^2\,\dd y\) in dimension four.
The coordinate functions of the round region have Rayleigh quotients
near the round value in a representative that is round there. In the
minimizing representative every function carried by the round region
has Rayleigh quotient of order at least \(\rho^{-2}\), on a set of
volume \(O(\rho^4)\). Trial quotients computed in a glued
representative give no information about \(3\lambda_5/s\) on the
minimizer.

A change of the conformal class inside a small ball cannot raise the
Yamabe constant appreciably, whatever is placed in the ball.

\begin{proposition}\label{prop:small-ball}
Let \(g\) be a volume-one Yamabe minimizer on a closed four-manifold
\(M\), and let \(p\in M\). There are constants \(C\) and \(r_0>0\),
depending on \(g\), with the following property for \(0<r\le r_0\).
Let \((M',g')\) be a closed Riemannian four-manifold, and let \(\Phi\)
be a diffeomorphism from \(M\setminus\overline{B_r(p)}\) onto the
complement of a compact subset of \(M'\) such that \(\Phi^*g'\) is
conformal to \(g\). Then
\[
 Y([g'])\le Y([g])+\frac{C\,r}{\log^2(1/r)} .
\]
\end{proposition}

\begin{proof}
Let \(d\) be the \(g\)-distance to \(p\), put \(R=\sqrt r\), and let
\(\chi\) equal \(0\) on \(B_{2r}(p)\), \(1\) outside \(B_R(p)\), and
\(\log(d/2r)/\log(R/2r)\) in between. For small \(R\) the geodesic
spheres of radius \(d\le R\) have volume at most \(Cd^3\), so the
coarea formula gives
\[
 \int_M\abs{\dd\chi}^2\,\dv
 \le\frac{C}{\log^2(R/2r)}\int_{2r}^R d\,\dd d
 \le\frac{C'\,r}{\log^2(1/r)},\qquad
 \int_M\chi^4\,\dv\ge1-C'r^2 .
\]
Since \(s_g=Y([g])\) is constant and \(1-C'r^2\le\int\chi^2\,\dv\le1\),
the Yamabe quotient of \(\chi\) for \(g\) is at most
\(Y([g])+C''r/\log^2(1/r)\), for either sign of \(Y([g])\). The
support of \(\chi\) is a compact subset of the domain of \(\Phi\). If
\(\Phi^*g'=a^2g\) there, the function equal to
\((\chi/a)\circ\Phi^{-1}\) on the image of \(\Phi\) and to zero
elsewhere has the same Yamabe quotient for \(g'\), by the covariance
\eqref{eq:conformal-L}. It is an admissible test function for
\(Y([g'])\).
\end{proof}

Consequently a sequence of conformal classes that agree with one fixed
class \([g_0]\) outside balls shrinking to a point has
\(\limsup Y\le Y([g_0])\), however the balls are filled, and
\(Y([g_0])<Y(\Sphere)\) unless \([g_0]\) is the round
class~\cite{LeeParker1987Yamabe,Schoen1984Yamabe}. The gluing of
Proposition~\ref{prop:flat-neck} is the case in which nothing changes
at all.

\subsection{Delaunay minimizers on \texorpdfstring{\(S^1\times S^3\)}{S1 x S3}}
\label{sec:necklace}

Let \(g_0=\dd t^2+g_{S^3}\) on \(S^1(\ell)\times S^3\), where the circle
has length \(\ell\) and \(g_{S^3}\) is the unit round metric. Its
scalar curvature is \(6\). For a positive function \(u\) of \(t\),
\eqref{eq:conformal-L} without the map potential gives
\begin{equation}\label{eq:cyl-scalar}
 s_{u^2g_0}=u^{-3}(-6u''+6u).
\end{equation}
Fix \(m\in(0,1)\) and write \(m_1=1-m\). Let \(\sn\), \(\cn\), \(\dn\)
be the Jacobian elliptic functions of modulus \(\sqrt m\), and let
\(K\), \(E\) be the complete elliptic integrals of the first and second
kind of the same modulus. The function \(\dn\) is even, has minimal
period \(2K\), and decreases from \(1\) at \(0\) to \(\sqrt{m_1}\) at
\(K\). It satisfies
\begin{equation}\label{eq:dn-ode}
 (\dn')^2=(1-\dn^2)(\dn^2-m_1),\qquad
 \dn''=(2-m)\dn-2\dn^3 ,
\end{equation}
by \cite[(22.13.3)]{DLMF} and differentiation. Put
\begin{equation}\label{eq:delaunay}
 u_m(t)=\dn\Bigl(\frac{t}{\sqrt{2-m}}\Bigr),\qquad
 g_m=u_m^2g_0,\qquad T(m)=2K\sqrt{2-m}.
\end{equation}
Then \(u_m''=u_m-2u_m^3/(2-m)\). By \eqref{eq:cyl-scalar} the metric
\(g_m\) has constant scalar curvature \(s=12/(2-m)\) on every circle
whose length is a multiple of the minimal period \(T(m)\) of \(u_m\).
In the arclength \(r\) along the axis, \(g_m=\dd r^2+u_m^2g_{S^3}\) is
a periodic chain of bulges with equators of radius \(1\), joined by
necks of radius \(\sqrt{m_1}\). As \(m\to1\) the profile tends to
\(1/\cosh t\), the unit round sphere in cylindrical coordinates. As
\(m\to0\) it tends to the product metric.

Let \(I_{2j}=\int_0^K\dn^{2j}\,\dd x\). The substitution
\(\sn x=\sin\varphi\) gives \(I_0=K\) and \(I_2=E\). Integrating the
derivative of \(\dn^{2j-1}\dn'\) over \([0,K]\) and using
\eqref{eq:dn-ode} gives
\((2j+1)I_{2j+2}=2j(2-m)I_{2j}-(2j-1)m_1I_{2j-2}\). Hence
\begin{align}
 I_4&=\tfrac13\bigl(2(2-m)E-m_1K\bigr),\qquad
 I_6=\tfrac15\bigl(4(2-m)I_4-3m_1E\bigr),\notag\\
 R(m)&=\frac{4(2-m)I_4-I_6-m_1E}{4I_6}.
 \label{eq:I4I6}
\end{align}

\begin{theorem}\label{thm:necklace}
The function \(T\) is an increasing bijection from \((0,1)\) onto
\((\sqrt2\,\pi,\infty)\). Let \(\ell>\sqrt2\,\pi\) and let \(m\) be the
solution of \(T(m)=\ell\). On \(S^1(\ell)\times S^3\) the following
hold.
\begin{enumerate}
\item[(i)] The Yamabe minimizers of \([g_0]\) are exactly the metrics
\(c\,u_m(t-t_0)^2g_0\) with \(c>0\) and \(t_0\in\RR\). They are
homothetic to \(g_m\), and
\[
 Y([g_0])=24\pi\,(2-m)^{-3/4}I_4^{1/2}<Y(\Sphere).
\]
\item[(ii)] Every minimizer satisfies \(s=3\lambda_1\), and
\(u_m'/u_m\) is a first eigenfunction of \(g_m\).
\item[(iii)] Every minimizer \(g\) satisfies
\[
 0\le\frac{3\lambda_5}{s}-1\le R(m)-1,\qquad
 \frac{A(g)}{s}\le\tfrac45\bigl(R(m)-1\bigr).
\]
\item[(iv)] Every minimizer \(g\) satisfies
\(\int\abs{\Ric_0}^2\,\dv=s^2\Vol(g)/12\), where \(\Ric_0\) is the
trace-free Ricci tensor.
\item[(v)] As \(\ell\to\infty\), with \(\Lambda=\log(16/m_1)\),
\[
 \ell=\Lambda+O(m_1\Lambda),\qquad
 R(m)-1=\tfrac{15}{16}m_1+O(m_1^2\Lambda),\qquad
 1-\frac{Y([g_0])}{Y(\Sphere)}=\tfrac38m_1+O(m_1^2\Lambda).
\]
\end{enumerate}
In particular the volume-one minimizers \(g_\ell\) of
\(S^1(\ell)\times S^3\) have \(\e(g_\ell)\to0\) and
\(A(g_\ell)/s\to0\) as \(\ell\to\infty\), with
\(\e(g_\ell)\le(\tfrac52+o(1))\bigl(1-Y([g_0])/Y(\Sphere)\bigr)\).
For \(\ell\le\sqrt2\,\pi\) the minimizers are the multiples of \(g_0\).
\end{theorem}

\begin{proof}
\emph{Reduction to an ordinary differential equation.}
The class \([g_0]\) contains a metric of positive scalar curvature, so
a metric \(u^2g_0\) of constant scalar curvature \(s\) in it has
\(s>0\). Lift \(u\) to \(\RR\times S^3\) and identify this cylinder
with \(\RR^4\setminus\{0\}\) by \(y=e^t\theta\), so that
\(\abs{\dd y}^2=e^{2t}g_0\). By \eqref{eq:conformal-L} the function
\(v=u/\abs y\) satisfies \(6\Delta v=sv^3\) for the flat Laplacian.
It is unbounded at the origin because \(u\) is periodic and positive.
The symmetry theorem of Caffarelli, Gidas, and
Spruck~\cite[Theorem~8.1]{CaffarelliGidasSpruck1989}, applied with the
nonlinearity \(sv^3/6\), shows that \(v\) is radially symmetric about
the origin. Thus \(u\) is a positive \(\ell\)-periodic function of
\(t\) with \(-6u''+6u=su^3\). This reduction is due to
Schoen~\cite[\S2]{Schoen1989Variational}.

\emph{Periodic solutions.}
The quantity \(u'^2-u^2+su^4/12\) is constant along solutions. Let
\(u\) be a nonconstant positive periodic solution. After a translation
and multiplication of \(u\) by a constant, which rescales \(s\), its
maximum is \(u(0)=1\). Equating the first integral at the maximum and
at the minimum \(u_{\min}\in(0,1)\) gives \(u_{\min}^2=12/s-1\). Hence
\(6<s<12\), and \(s=12/(2-m)\) for a unique \(m\in(0,1)\). By
\eqref{eq:dn-ode}, \(u_m\) solves the same initial value problem, so
\(u=u_m\). The metrics of constant scalar curvature in \([g_0]\) are
therefore the multiples of \(g_0\) and of the metrics
\(u_m(t-t_0)^2g_0\) with \(T(m)=\ell/j\) for an integer \(j\ge1\).

\emph{The period.}
By \cite[(19.4.1), (19.4.2)]{DLMF}, written in the variable \(m\),
\(2mm_1\,\dd K/\dd m=E-m_1K\) and \(2m\,\dd E/\dd m=E-K\). Therefore
\[
 T'(m)=\frac{F(m)}{mm_1\sqrt{2-m}},\qquad
 F(m)=(2-m)E-2m_1K,\qquad F'(m)=\tfrac32(K-E),
\]
and \(F(0)=0\). The integrand of \(K\) exceeds that of \(E\), so
\(F'>0\) and \(T\) is increasing. Its limit at \(m=0\) is
\(\sqrt2\,\pi\), because \(K=\pi/2\) there, and its limit at \(m=1\)
is infinite. The count of solutions in
\cite[\S2]{Schoen1989Variational} depends on this monotonicity.

\emph{Stability.}
A Yamabe minimizer exists and is a smooth metric of constant scalar
curvature~\cite{LeeParker1987Yamabe,Schoen1984Yamabe}. It satisfies
\(s\le3\lambda_1\) by the second variation used for
\eqref{eq:balance-chain}. The product \(g_0\) has \(s=6\) and
\(\lambda_1\le(2\pi/\ell)^2<2\) when \(\ell>\sqrt2\,\pi\), so it is not
a minimizer. Let \(u=u_m\) with \(T(m)=\ell/j\). For a function
\(\phi\) of \(t\) alone,
\[
 \Delta_{u^2g_0}\phi=-u^{-4}(u^2\phi')',\qquad
 \frac{\int\abs{\dd\phi}^2\,\dd V_{u^2g_0}}{\int\phi^2\,\dd V_{u^2g_0}}
 =\frac{\int u^2\phi'^2\,\dd t}{\int u^4\phi^2\,\dd t}.
\]
The field \(\partial_t\) is conformal for \(u^2g_0\), with
\(\mathcal L_{\partial_t}(u^2g_0)=2\sigma\,u^2g_0\) and
\(\sigma=u'/u\). By \eqref{eq:conformal-eigen},
\begin{equation}\label{eq:sigma-mode}
 -(u^2\sigma')'=\tfrac{s}{3}u^4\sigma .
\end{equation}
The zeros of \(\sigma\) on the circle are the \(2j\) critical points
of \(u\). Suppose \(j\ge2\). Let \(J_1,J_2,J_3\) be three of the
\(2j\) arcs between consecutive zeros, and let \(\sigma_i\) equal
\(\sigma\) on \(J_i\) and zero elsewhere. Multiplying
\eqref{eq:sigma-mode} by \(\sigma\) and integrating over \(J_i\) gives
\(\int u^2\sigma_i'^2\,\dd t=(s/3)\int u^4\sigma_i^2\,\dd t\).
Choose \((a_1,a_2)\ne0\) so that \(\psi=a_1\sigma_1+a_2\sigma_2\)
satisfies \(\int u^4\psi\,\dd t=0\). Then \(\psi\) is a nonzero
Lipschitz function of mean zero with Rayleigh quotient \(s/3\). If
\(\lambda_1\ge s/3\), it would minimize the Rayleigh quotient among
functions of mean zero and so would be a first eigenfunction. A
solution of the linear equation \eqref{eq:sigma-mode} that vanishes on
the arc \(J_3\) is zero, a contradiction. Hence \(3\lambda_1<s\) for
\(j\ge2\), and these metrics are not minimizers. The minimizers are
therefore the metrics with \(j=1\), which are translates and multiples
of \(g_m\) with \(T(m)=\ell\) and are isometric up to scale. For
\(\ell\le\sqrt2\,\pi\) no \(m\) and \(j\) satisfy \(T(m)=\ell/j\), and
only the multiples of \(g_0\) remain.

\emph{Energy and the first eigenvalue.}
The volume of \(g_m\) on the circle of length \(T(m)\) is
\(2\pi^2\int_0^{T(m)}u_m^4\,\dd t=4\pi^2\sqrt{2-m}\,I_4\), and
\(Y([g_0])=s\Vol(g_m)^{1/2}\), which is the formula in (i). The
strict inequality holds because the class is not that of the round
sphere~\cite{LeeParker1987Yamabe,Schoen1984Yamabe}. For (ii), the
minimizer has \(s\le3\lambda_1\), and \(\sigma\) is a nonconstant
eigenfunction with eigenvalue \(s/3\) by \eqref{eq:sigma-mode}.

\emph{The transverse trial functions.}
Let \(\theta_1,\ldots,\theta_4\) be the restrictions to \(S^3\) of the
coordinates of \(\RR^4\). They satisfy
\(\Delta_{S^3}\theta_i=3\theta_i\), and \(1,\theta_1,\ldots,\theta_4\)
are orthogonal in \(L^2(S^3)\). For \(\phi=f(t)\theta_i\) and
\(g=u^2g_0\),
\[
 \frac{\int\abs{\dd\phi}_g^2\,\dv}{\int\phi^2\,\dv}
 =\frac{\int u^2(f'^2+3f^2)\,\dd t}{\int u^4f^2\,\dd t}.
\]
Take \(u=f=u_m\) and \(x=t/\sqrt{2-m}\). By \eqref{eq:dn-ode},
\(\int_0^K\dn^2(\dn')^2\,\dd x=(2-m)I_4-I_6-m_1E\), and the quotient
equals \((s/3)R(m)\). The five functions \(\sigma\) and
\(u_m\theta_i\) have mean zero and are mutually orthogonal, both in
\(L^2(\dv)\) and in the Dirichlet form, by the orthogonality on
\(S^3\). On the span of the four functions \(u_m\theta_i\) the Rayleigh
quotient is identically \((s/3)R(m)\). The min--max principle gives
\(\lambda_4\le(s/3)R(m)\), and so \(R(m)\ge1\) by (ii). On the span of
all five functions the Rayleigh quotient is then at most
\((s/3)R(m)\), which gives \(\lambda_5\le(s/3)R(m)\). The bound for
\(A(g)\) follows from
\(A(g)/s=\tfrac15\sum_{i\le5}(3\lambda_i/s-1)\) and (ii).

\emph{Curvature and asymptotics.}
The metric \(g_m\) is locally conformally flat and
\(\chi(S^1\times S^3)=0\), so (iv) is the Chern--Gauss--Bonnet formula
\eqref{eq:gauss-bonnet} below. For (v), the expansions
\cite[(19.12.1), (19.12.2)]{DLMF} give
\(K=\tfrac12\Lambda+O(m_1\Lambda)\) and
\(E=1+\tfrac14m_1(\Lambda-1)+O(m_1^2\Lambda)\). The logarithms cancel
at first order in
\[
 I_4=\tfrac23+\tfrac12m_1+O(m_1^2\Lambda),\qquad
 I_6=\tfrac8{15}+\tfrac13m_1+O(m_1^2\Lambda).
\]
Then \eqref{eq:I4I6} gives the expansion of \(R(m)\), and (i) gives
\(Y([g_0])^2=384\pi^2(1-\tfrac34m_1)+O(m_1^2\Lambda)\). The first
relation in (v) is \(\ell=T(m)\). The final assertions follow from
(ii), (iii), and (v).
\end{proof}

Theorem~\ref{thm:necklace} shows that relative near equality in
Theorem~\ref{thm:spectral} occurs on a closed four-manifold with
infinite fundamental group. The hypotheses of Theorem~\ref{thm:near}
hold along \(g_\ell\), so the joint law of the five lowest modes
converges to the round law while the manifold remains
\(S^1\times S^3\). The supremum of the Yamabe constants of the classes
\([g_0]\) is \(Y(\Sphere)\) and is not attained, in agreement with
Schoen's description~\cite[\S2]{Schoen1989Variational}.

\paragraph{Separation of variables and a numerical complement.}
The eigenvalue problem of \(g_m\) separates exactly. For
\(\phi=f(x)Y(\theta)\), with \(x=t/\sqrt{2-m}\) and \(Y\) a spherical
harmonic of degree \(n\) on \(S^3\), the function \(h=\dn\cdot f\) and
the scaled eigenvalue \(\nu=3\lambda/s\) of Theorem~\ref{thm:near}
satisfy the Lam\'e equation
\begin{equation}\label{eq:lame}
 -h''+(2-m)(n+1)^2h=(4\nu+2)\dn^2h\qquad\text{on }\RR/2K\mathbb Z .
\end{equation}
Three solutions are elementary. They are \(h=\dn\) with \(n=0\) and
\(\nu=0\), the constants; \(h=\sn\cn\) with \(n=0\) and \(\nu=1\), the
mode \(\sigma\); and \(h=\sn\cn\dn\) with \(n=1\) and
\(\nu=\tfrac52\). The remaining statements in this paragraph and the
entries of Table~\ref{tab:necklace} are double-precision computations.
They are not used in any proof. A Fourier--Galerkin discretization of
\eqref{eq:lame}, checked against a finite-difference scheme, was run on
93 values of \(m\) between \(0.002\) and \(1-10^{-7}\). It indicates
that the other eigenvalues below \(\tfrac52\) are a simple eigenvalue
\(\nu_e(m)\) in the sector \(n=0\), which increases from \(1\) to
\(\tfrac52\) with \(m\), and the ground state \(\nu_1(m)\) of the
sector \(n=1\), of multiplicity four, which decreases from
\(\tfrac32\) to \(1\). In either order of these two,
\(\e(g_m)=\nu_1(m)-1\). The computed values satisfy
\((\nu_1-1)/m_1=0.9375-0.387\,m_1+\cdots\) for small \(m_1\). The bound
in (iii) is thus attained to first order, and
\(\e\approx\tfrac52(1-Y([g_0])/Y(\Sphere))\) for a thin neck.

\begin{table}[ht]
\centering\small
\setlength{\tabcolsep}{5pt}
\begin{tabular}{ccccccccc}
\hline
\(m_1\) & \(\ell\) & \(\e\) & \(\e/m_1\) & \(D/m_1\)
 & \(\max\abs{q-\tfrac{15}{7}}\) & \(\norm{\abs X^2-1}_2\)
 & \(\inf\abs X\) & \(\int e^2\)\\
\hline
\(10^{-1}\) & \(5.4079\) & \(9.0072\times10^{-2}\) & \(0.9007\) & \(0.2962\)
 & \(2.7\times10^{-1}\) & \(1.2\times10^{-1}\) & \(0.154\) & \(1.1690\)\\
\(10^{-2}\) & \(7.4281\) & \(9.3367\times10^{-3}\) & \(0.9337\) & \(0.3633\)
 & \(3.9\times10^{-2}\) & \(1.8\times10^{-2}\) & \(0.0494\) & \(1.0535\)\\
\(10^{-4}\) & \(11.9838\) & \(9.3746\times10^{-5}\) & \(0.93746\) & \(0.37482\)
 & \(4.3\times10^{-4}\) & \(2.6\times10^{-4}\) & \(0.0050\) & \(1.0430\)\\
\(10^{-6}\) & \(16.5881\) & \(9.3750\times10^{-7}\) & \(0.937500\) & \(0.374998\)
 & \(4.3\times10^{-6}\) & \(3.2\times10^{-6}\) & \(0.0005\) & \(1.0429\)\\
\hline
\end{tabular}
\caption{Computed quantities for the volume-one Delaunay minimizers.
Here \(D=1-Y([g_0])/Y(\Sphere)\), \(X\) is the map of
Theorem~\ref{thm:near}, \(q=\int(\sum_ia_iu_i)^4\,\dv\) is evaluated
for unit vectors \(a\) in the axial direction, a transverse direction,
and their diagonal, and \(e=\sum_i\Gamma(X_i,X_i)\).}
\label{tab:necklace}
\end{table}

The fourth moments and the radial error converge as
Theorem~\ref{thm:near} requires, with room in its rates. Three
quantities do not converge to their round values. First, the ground
state of the sector \(n=1\) is positive, so \(X\) has no zeros and
\(X/\abs X\) omits the two poles of the axis. This map from
\(S^1\times S^3\) to \(S^4\) has degree zero, and
\(\inf\abs X\approx\tfrac12\sqrt{m_1}\). The bulge covers the sphere
once and the neck returns along a thin tube around the axis. Second,
the energy density \(e\), which is identically \(1\) for the round
sphere, has \(\int e^2\,\dv\to1.043\). The Gram defect
\(\Gamma(X_1,X_1)-\tfrac14(1-X_1^2)\) of the axial coordinate tends to
zero in \(L^1\), while its \(L^2\) norm tends to \(0.207\). Third, the
curvature does not converge. In the notation of \eqref{eq:delaunay},
a direct calculation gives
\begin{equation}\label{eq:necklace-ricci}
 \Ric_0=\frac{c}{u_m^2}\bigl(-3\,\dd t^2+g_{S^3}\bigr),\qquad
 \abs{\Ric_0}^2=\frac{12c^2}{u_m^8},\qquad c=\frac{m_1}{2-m},
\end{equation}
which is exact. The least Ricci eigenvalue is attained at the necks
and equals \((s/4)(1-1/m_1)\). It is negative for every \(m\), so the
Ricci hypothesis of Aubry's theorem fails there, and the mass of
\(\abs{\Ric_0}^2\) in (iv) sits in the necks as \(m_1\to0\). Thus
convergence of the law of \(X\) does not
control the degree of \(X/\abs X\), the \(L^2\) norm of the Gram
defect, or the curvature. Each is affected by a region whose volume
tends to zero.

\subsection{An integral Ricci criterion at a Yamabe metric}
\label{sec:budget}

Write \(\Ric_0=\Ric-(s_g/4)g\) for the trace-free Ricci tensor. On a
closed four-manifold the Chern--Gauss--Bonnet formula is
\begin{equation}\label{eq:gauss-bonnet}
 8\pi^2\chi(M)=\int_M\Bigl(\tfrac14\abs{W}_{4\mathrm{ind}}^2
   -\tfrac12\abs{\Ric_0}^2+\tfrac1{24}s_g^2\Bigr)\,\dv ,
\end{equation}
as in \cite[Eq.~(1.2)]{ChangQingYang2006Gap}. On an oriented manifold,
\(\tfrac14\int\abs{W}_{4\mathrm{ind}}^2\,\dv=E_++E_-\) and
\eqref{eq:signature} give
\begin{equation}\label{eq:eplus}
 E_+(g)=2\pi^2\bigl(2\chi(M)+3\tau(M)\bigr)
   +\tfrac14\int_M\abs{\Ric_0}^2\,\dv-\tfrac1{48}\int_Ms_g^2\,\dv .
\end{equation}

\begin{proposition}\label{prop:budget}
Let \(g\) be a volume-one metric of constant scalar curvature \(s>0\)
on a closed four-manifold \(M\). If
\begin{equation}\label{eq:ricci-budget}
 \int_M\abs{\Ric_0}^2\,\dv<\frac{s^2}{12}-8\pi^2\chi(M),
\end{equation}
then \(M\) is diffeomorphic to \(\Sphere\) or \(\mathbb{RP}^4\). If
\(M\) is a smooth homotopy four-sphere, the condition reads
\(\int_M\abs{\Ric_0}^2\,\dv<s^2/12-16\pi^2\), the conclusion is that
\(M\) is diffeomorphic to \(\Sphere\), and
\begin{equation}\label{eq:eplus-sphere}
 E_+(g)=\frac{Y(\Sphere)^2-s^2}{48}
        +\tfrac14\int_M\abs{\Ric_0}^2\,\dv .
\end{equation}
\end{proposition}

\begin{proof}
A metric of positive scalar curvature lies in a class of positive
Yamabe constant. By \eqref{eq:gauss-bonnet}, the inequality
\eqref{eq:ricci-budget} is equivalent to
\(\int\abs{W}_{4\mathrm{ind}}^2\,\dv<16\pi^2\chi(M)\), which is the
hypothesis of~\cite[Theorem~A]{ChangGurskyYang2003}. On a homotopy
sphere \(\chi=2\) and \(\tau=0\), simple connectivity excludes
\(\mathbb{RP}^4\), and \eqref{eq:eplus} with
\(Y(\Sphere)^2=384\pi^2\) is \eqref{eq:eplus-sphere}.
\end{proof}

With \(s^2\) replaced by \(\int s_g^2\,\dv\), the difference of the two
sides of \eqref{eq:ricci-budget} is a conformal invariant, so the
criterion may be tested on any representative. A Yamabe minimizer
gives it the form above with \(s=\Y\). On a homotopy four-sphere the
threshold \(s^2/12-16\pi^2\) is then positive exactly when
\(\Y>8\sqrt3\,\pi=Y(\Sphere)/\sqrt2\), and it increases to
\(16\pi^2\) as \(\Y\to Y(\Sphere)\). By \eqref{eq:eplus-sphere}, every
class on a homotopy four-sphere has
\(E_+\ge(Y(\Sphere)^2-\Y^2)/48\), with equality exactly when a
minimizer is Einstein. This is the form taken here by the bound of
\cite[Eq.~(1.7)]{ChangQingYang2006Gap}. On a homotopy four-sphere that
is not diffeomorphic to \(\Sphere\), every volume-one Yamabe minimizer
with \(s>0\) would satisfy
\(\int\abs{\Ric_0}^2\,\dv\ge s^2/12-16\pi^2\).

When equality holds in the hypothesis of Chang, Gursky, and Yang and
\(M\) is neither \(\Sphere\) nor \(\mathbb{RP}^4\), the class is that
of the Fubini--Study metric or of a quotient of the product
\(S^3\times S^1\)~\cite[Theorem~1.1]{ChangQingYang2006Gap}. The
Delaunay metrics are in this equality case, since \(\chi=0\) and
\(\int\abs{\Ric_0}^2\,\dv=s^2\Vol/12\) by Theorem~\ref{thm:necklace}.
The family with \(\e\to0\) therefore misses \eqref{eq:ricci-budget} by
exact equality at every length.

A pointwise Ricci bound implies the integral bound without spectral
input. Suppose \(\Ric_g\ge(1-\delta)(s/4)g\) for a volume-one metric
of constant scalar curvature \(s>0\). The eigenvalues of \(\Ric_0\)
are at least \(-\delta s/4\) and sum to zero, so
\(\abs{\Ric_0}^2\le\tfrac34\delta^2s^2\), the extreme case being
\((3,-1,-1,-1)\delta s/4\). Hence \eqref{eq:ricci-budget} holds when
\[
 \delta^2<\frac19-\frac{32\pi^2\chi(M)}{3s^2}.
\]
On a homotopy four-sphere this is
\(\delta^2<\tfrac19-64\pi^2/(3s^2)\), which allows \(\delta\) up to
\(1/(3\sqrt2)\approx0.236\) as \(s\to Y(\Sphere)\). The route through
Aubry's theorem in Section~\ref{sec:near} needs both \(\delta\) and
\(\e\) small and gives a Gromov--Hausdorff estimate as well. The
integral criterion uses the Euler characteristic, and on a homotopy
four-sphere it needs \(s>8\sqrt3\,\pi\).

Relative pinching does not imply \eqref{eq:ricci-budget}, and the
Bochner formula shows why. Let \(g\) have constant scalar curvature
\(s\), let \(\Delta_g\phi=\lambda\phi\), and write
\(\nabla^2\phi=\Hess_0\phi-\tfrac14(\Delta_g\phi)g\) and
\(\nu=3\lambda/s\). The integrated Bochner formula
\(\int\abs{\nabla^2\phi}^2\,\dv+\int\Ric(\nabla\phi,\nabla\phi)\,\dv
=\int(\Delta_g\phi)^2\,\dv\) gives
\begin{equation}\label{eq:bochner-pinch}
 \int_M\abs{\Hess_0\phi}^2\,\dv
 +\int_M\Ric_0(\nabla\phi,\nabla\phi)\,\dv
 =\frac{s^2}{12}\,\nu(\nu-1)\int_M\phi^2\,\dv .
\end{equation}
For the first five modes of a volume-one Yamabe minimizer, normalized
in \(L^2\), the right side lies between \(0\) and
\(\tfrac{s^2}{12}\e(1+\e)\). Relative pinching thus controls the sum
of the two terms on the left. The first is nonnegative and the second
has no sign. For the mode \(\sigma=u_m'/u_m\) of the Delaunay metric,
\(\nu=1\) and a direct calculation gives
\(\Hess_0\sigma=\sigma\Ric_0\). The function \(\sqrt{m_1}/\dn\) solves
the second equation in \eqref{eq:dn-ode} with the initial data of
\(\dn(\cdot+K)\), so \(u_m(t+T(m)/2)=\sqrt{m_1}/u_m(t)\). With
\eqref{eq:necklace-ricci} this symmetry gives
\[
 \int\abs{\Hess_0\sigma}^2\,\dv
 =-\int\Ric_0(\nabla\sigma,\nabla\sigma)\,\dv
 =\int\sigma^2\abs{\Ric_0}^2\,\dv=\frac{s^2}{12}\int\sigma^2\,\dv
\]
at every length. A pinched mode can therefore have a trace-free
Hessian of fixed size in \(L^2\), carried by the region where
\(\Ric_0\) is large.

\section{Dirac transmission and curvature on marked collars}\label{sec:collar-transmission}

A comparison collar gives a way to use the canonical fields of
Section~\ref{sec:green} on the actual closed source. There are two
analytic requirements: solving the transmission problem for the
prescribed traces and controlling the projected derivative in
\eqref{eq:weyl-defect}. We give the estimates used at this interface.
The strip analysis, Cauchy data from actual capped regions and
quantitative comparison of end fields are developed in
\cite{Baker2026DiracTransmission}. The geometry that supplies or tests
the end identifications is treated in
\cite{Baker2026MarkedNecks,Baker2026MarkedLimits}.

\subsection{A global correction for seam jumps}

Let \((M,g)\) be a smooth closed oriented Riemannian spin four-manifold
with \(Y([g])\ge y_0>0\). All Dirac operators in this section are
untwisted. Supply an embedded oriented collar with coordinates
\([-2,L+2]\times S^3\), \(L\ge6\), and a smooth conformal representative
\(h=F^2g\) in those coordinates, where \(F>0\). Near \(t_0=0\) and
\(t_1=L\), on the two windows \(|t-t_\nu|\le3/2\), suppose
\(h\) is close to \(h_0=dt^2+g_{S^3(1)}\) in the norm
\(\|h-h_0\|_{L^\infty}+\|\nabla^0(h-h_0)\|_{L^4}\).
Use the same symmetric spin-frame identification on both sides of
each cut. Let \(A\) be the reference tangential Dirac operator with
normal \(+\partial_t\), and equip the trace space
\(\mathcal T=H^{1/2}(S^3,S^\pm_0)\) with the equivalent norm
\(\||A|^{1/2}\cdot\|_2\). The reference spectrum has
\(\abs{A}\ge3/2\)~\cite{Bar1996SphereDirac}. Write
\(J_\nu q\) for the positive-side trace minus the negative-side trace
in the reference frame after harmonic conformal transport
\eqref{eq:conformal-spin}.

\begin{proposition}[Closed-source jump correction]\label{prop:dirac-jumps}
For fixed chirality and prescribed jumps \(j_\nu\in\mathcal T\),
there is a unique piecewise \(H^1\) Dirac-harmonic spinor \(q\) with
\(J_\nu q=j_\nu\) on the cut-open source. If the two normalized
windows meet the stated uniform smallness bound, then
\begin{equation}\label{eq:dirac-jumps}
 \sum_{\nu=0}^1\sum_{\sigma=\pm}
 \left(\norm{q_h}_{H^1(B_\nu^\sigma)}^2+
             \norm{\tau_{\nu,\sigma}q}_{\mathcal T}^2\right)
 \le C(y_0)^2\sum_{\nu=0}^1\norm{j_\nu}_{\mathcal T}^2 ,
\end{equation}
where \(B_\nu^\sigma\) are the half-windows of width \(1/2\) at each
cut, and \(\tau_{\nu,\sigma}\) denotes the corresponding trace.
The global critical norm also satisfies
\(\|q\|_{L^{8/3}(g)}\le C(y_0)
(\sum_\nu\|j_\nu\|_{\mathcal T}^2)^{1/2}\).
The constants do not depend on \(L\), the original conformal scales
at the seams or the geometry outside the windows.
\end{proposition}

\begin{proof}
Extend the normalized representative to \(\widetilde g=b^2g\)
equal to \(h\) on both unit windows. A smooth positive \(b\) can be
chosen by interpolating \(F\) to a sufficiently small positive
constant outside the disjoint buffered windows. Convexity of
\(x\mapsto x^4\) bounds the total \(\widetilde g\)-volume by the sum
of the two buffered metric volumes and one. On any conformal
representative \(k\) in the closed class, the Hijazi estimate yields
\[
 \|D_k^{-1}\|_{L^2(k)\to L^2(k)}
 \le (3/y_0)^{1/2}\Vol(k)^{1/4}
\]
~\cite[Theorem A and (2.1)--(2.2)]{Hijazi1986Conformal}.
In these conventions the first eigenvalue \(\mu_1\) of
\(6\Delta_k+s_k\) is at least \(Y([g])\Vol(k)^{-1/2}\).
Its positive first eigenfunction provides a conformal gauge in
which the integrated twistor identity \eqref{eq:energy} gives
\(3\lambda^2\ge\mu_1\) for every closed Dirac eigenvalue \(\lambda\).

For \(s=t-t_\nu\ge0\) lift each jump to the positive half-window
as \(v_\nu(s)=\zeta(s)e^{-s|A|}j_\nu\), transported into the normalized
frame. The cutoff \(\zeta\) equals one near zero and vanishes for
\(s\ge1/2\). Integration of \((1+\mu^2)e^{-2\mu s}\) on each
absolute angular mode \(\mu\ge3/2\) gives
\(\|v_\nu\|_{H^1}\le C\|j_\nu\|_{\mathcal T}\). Extend by zero on
the other sides, set \(v=v_0+v_1\), and let \(G=D_{\widetilde
g,\mathrm{pw}}v\) be the regular Dirac image on the open pieces.
The first-order operator comparison bounds \(\|G\|_2\) by the
sum of the lift norms. Solving on the closed source gives
\(w=D_{\widetilde g}^{-1}G\). This corrects the regular forcing,
while the seam delta terms of the distributional derivative of
\(v\) encode the jumps. The field \(v-w\) is harmonic on the
pieces and has the prescribed jumps.

In the windows the operator difference from the round product is
bounded by \(C(\|h-h_0\|_\infty+\|\nabla^0(h-h_0)\|_4)\) as a map
from \(H^1\) to \(L^2\). The zero-order term with one metric
derivative uses \(H^1\subset L^4\). The product interior estimate
with a cutoff absorbs this error and gives the uniform \(H^1\)
bound for \(w\) on the half-windows from its global \(L^2\) bound
and \(\|G\|_2\). The trace estimate gives \eqref{eq:dirac-jumps}.
A field with zero jumps glues to a weak harmonic spinor on the
closed manifold and therefore vanishes by the closed spectral bound.

For the global estimate put \(p=8/5\), \(p'=8/3\).
Given smooth \(f,k\) with \(L^p(g)\) norms one, choose
\(a_\epsilon^4=(|f|_g^2+\epsilon)^{4/5}
+(|k|_g^2+\epsilon)^{4/5}\) and transport the two sources to
\(g_\epsilon=a_\epsilon^2g\) with weight \(-5/2\).
The pairing of \(D^{-1}f\) and \(k\) is conformally invariant
by \eqref{eq:conformal-spin}, and each transported \(L^2\) norm
is at most one because \(a_\epsilon^{-1}|f|^2\le |f|^{8/5}\),
and similarly for \(k\). The volume of \(g_\epsilon\) tends to
two. Applying the closed \(L^2\) bound in that gauge and taking
the dual supremum proves
\[
 \|D_g^{-1}f\|_{8/3,g}\le 2^{1/4}(3/y_0)^{1/2}\|f\|_{8/5,g}.
\]
Density extends this to \(L^{8/5}\) forcing. Apply the estimate
to \(G\), supported in fixed-volume normalized windows, and the
local Sobolev estimate to \(v\) there. The \(L^{8/3}\) solution
norm is invariant under Dirac transport. This proves the global
bound on \(q\).
\end{proof}

The estimate acts on actual source traces. To compare two specific
end fields one still needs both trace mismatches in a common
identification at their respective seams and control of the
regular forcing of the chosen interpolation. A small solution in
the critical mass norm does not bound its projected differentiated
twistor defect. The defect requires a different estimate.

\subsection{Curvature of actual normalized Green fields}

\begin{proposition}[Projected defect on a measured source region]\label{prop:actual-dirac-defect}
Let \((M,g)\) be a connected smooth closed oriented Riemannian spin
four-manifold with \(Y([g])>0\). For any point \(p\) and any nonzero
negative half-spinor \(\xi\) there, the positive half-spinor
\(\psi=D_g^{-1}(\delta_p\xi)\) is nontrivial and smooth harmonic on
\(M\setminus\{p\}\). Let \(Z\) be its interior zero set and put
\(\widehat\psi=\psi/|\psi|\) off \(Z\cup\{p\}\).
For every measurable \(U\subset M\),
\begin{equation}\label{eq:actual-dirac-defect}
 \int_{U\setminus(Z\cup\{p\})}
       \big|\Pi(d^\nabla P\widehat\psi)^+\big|^2\,dV_g
 =\tfrac12 E_+(g;U).
\end{equation}
The same density is obtained in every smooth positive conformal
representative on the compared open set when the harmonic field is
transported with its Dirac weight and then normalized.
\end{proposition}

\begin{proof}
Invertibility follows from positive Yamabe constant and the closed
Dirac bound above. Elliptic regularity gives a smooth solution away
from the pole, and chirality gives its positive sign. A nonzero
distribution supported at one point is a finite sum of derivatives
of a delta. The invertible principal Dirac symbol away from the
zero covector multiplies its highest-order coefficient polynomial
to a nonzero polynomial of one higher order. Thus a solution
supported at \(p\) cannot have order-zero Dirac image
\(\delta_p\xi\). The punctured manifold is connected, and B\"ar's
nodal-set theorem for harmonic Dirac spinors gives volume zero for
\(Z\)~\cite[Main Theorem]{Baer1997Nodal}. At every nonzero point,
apply Proposition~\ref{prop:weyl-defect} to \(\widehat\psi\)
and use \(|\widehat\psi|=1\). The norm density is half the chiral
Weyl density, which is conformally invariant in four dimensions.
\end{proof}

Consequently a Weyl-energy estimate on an actual residual region
controls the left-hand side of \eqref{eq:actual-dirac-defect} at any
spinor amplitude. In particular, suppose finitely many buffered
annular charts of the physical source, up to constant scale, have
metric perturbations \(v_j\) uniformly small in \(L^\infty\).
For cutoffs \(\eta_j\) that equal one on the protected regions and
use long log-radius transitions, varying the metric by the sum of
the pushed-forward \(\eta_j^4v_j\) tests their weighted chiral
energy against the total Bach norm. The first-derivative curvature
errors include \(\sum_j\|\eta_j Dv_j\|_4^4\). The conformally
invariant defect therefore tends to zero on the covered residual
along a sequence with uniformly bounded cover multiplicity and
physical \(L^2\) curvature, vanishing metric amplitudes and total
Bach norm, and vanishing summed bulk and transition errors.
To pass to the entire residual also requires vanishing uncovered
chiral energy~\cite{Baker2026TwistorSections,Baker2026DiracTransmission}.
Derivative bounds on individual unit slabs alone permit a bulk
sum that grows with the length. Even when the estimate applies,
zero sets need other control for arguments that use a global
spinor direction.

\subsection{Physical volume at the comparison faces}

A normalized collar frame can see very little of the measure of
the original Yamabe minimizer. Let \(g_i\) be volume-one positive
Yamabe minimizers with \(Y([g_i])\ge y_0\). If \(0\le\chi\le1\)
is one on a region of physical volume \(m\), its support has
physical volume at most \(s\), and its derivative is supported
on a transition of physical volume \(v\), then
\begin{equation}\label{eq:dirac-volume-witness}
 y_0(1-\sqrt{s})\sqrt m\ \le\
 6\int_M|d\chi|_{g_i}^2\,dV_{g_i}\ \le\
 6\left(\int_M|d\chi|_{g_i}^4\,dV_{g_i}\right)^{1/2}\sqrt v .
\end{equation}
For the first inequality use the Yamabe quotient and
\(\int\chi^2\le\sqrt{s}(\int\chi^4)^{1/2}\).
The second is H\"older on the transition. The fourth-power
gradient integral can be evaluated in the conformal comparison
gauge. Thus on an actual fixed marking with bounded cutoff
four-energy, two exteriors each carrying a positive fraction
of the physical volume give a positive compact volume witness.
If every fixed comparison compact evacuates, testing each exterior
separately gives \(\sqrt{\lambda_j}\le\lambda_j\) for every limiting
exterior volume \(\lambda_j\), so only one exit carries volume.

For the connection to smooth type, suppose the source is a fixed
oriented smooth homotopy four-sphere with vanishing total Bach
norm and \(4\pi^2\le E_+(g_i)\le8\pi^2-\delta\) for
\(\delta>0\). Supply a coherent comparison of whole compact
truncations of a smooth closed conformal base with finitely many
ordinary marks. The first-order comparison theorem uses
\(C^0\cap W^{1,4}\) convergence of the comparison metrics and a
local \(L^2\) scalar-curvature bound in that gauge. If the limiting
source Yamabe value is strictly subround, positive physical volume
on one fixed compared compact suffices for saturation.
At minimizing chiral energy over all smooth positive-Yamabe
classes on the source, a supplied nonround base gives the volume
saturation by itself and transfers the smooth type. The physical
sequence then converges in measured Gromov--Hausdorff distance
with no residual full-curvature energy
~\cite{Baker2026TwistorSections,Baker2026MarkedLimits}.
These implications help distinguish a physically relevant comparison
from one obtained by normalizing a tiny coordinate neighborhood.
Spectral convergence to a smooth limit for the weak comparison
used here is a separate condition for applying
Theorem~\ref{thm:spectral} to a pinched sequence.

\section{Outstanding challenges}\label{sec:open}

\paragraph{Construction on an arbitrary smooth homotopy sphere.}
Attaining \(A(g)=0\) at a local Yamabe minimizer would give the smooth
identification directly, by Proposition~\ref{prop:local}. A less rigid
construction would produce actual smooth volume-one Yamabe minimizers
with \(s_j>0\) and \(3\lambda_5(g_j)/s_j\to1\), together with the
geometric control discussed in the next paragraph.
Attaching a round region to a punctured core across a conformally flat
neck does not produce them. By Proposition~\ref{prop:flat-neck} the
conformal class and its minimizers are unchanged, and
\eqref{eq:chamber-rayleigh} places the round region at high frequency
and small volume in the minimizing representative. By
Proposition~\ref{prop:small-ball}, no change confined to a small ball
raises the Yamabe constant appreciably. A construction must therefore
change the conformal class on a region that carries a definite part of
the minimizer. On \(S^1\times S^3\) the relative gap and the Yamabe
deficit tend to zero together, by Theorem~\ref{thm:necklace}. Whether
\(\e\to0\) forces \(\Y\to Y(\Sphere)\) for volume-one Yamabe
minimizers, and whether the converse holds, are open questions. It is
also open whether every smooth homotopy four-sphere carries conformal
classes with Yamabe constant close to \(Y(\Sphere)\). Surgery gives
partial information. Positive scalar curvature survives surgery in
codimension at least three, by the theorem of Gromov and Lawson and of
Schoen and Yau as stated
in~\cite[Theorem~1.1]{BamlerLiMantoulidis2023Decomposing}. If
\(\sigma(M)\) denotes the supremum of the Yamabe constants of the
conformal classes on \(M\), surgery on a circle in a closed
four-manifold \(M\) gives a manifold \(N\) with
\(\sigma(N)\ge\min\{\sigma(M),\Lambda_{4,1}\}\) and
\(\Lambda_{4,1}>38.9\)~\cite[Theorem~1.2, Example~4.7]{AmmannDahlHumbert2015LowDim}.
A connected sum of copies of \(S^1\times S^3\) has
\(\sigma=Y(\Sphere)\), by Theorem~\ref{thm:necklace} and Kobayashi's
inequality for connected
sums~\cite{Kobayashi1987UnitVolume}\cite[Eq.~(14)]{LeBrun2021Scalar4}.
For a homotopy four-sphere obtained from such a sum by surgery on
circles, the resulting lower bound \(38.9\) is below the level
\(8\sqrt3\,\pi\approx43.5\) at which \eqref{eq:ricci-budget} can hold,
and well below \(Y(\Sphere)\approx61.6\).

\paragraph{Smooth recognition under approximation.}
Theorem~\ref{thm:near} gives the round law of five values under
relative pinching. The Delaunay metrics of
Theorem~\ref{thm:necklace} show that this law, the radial estimate,
and the weak Gram estimate are compatible with a handle of vanishing
volume. The estimates use no topological property of the manifold,
and near equality does not constrain one. A smooth conclusion needs
geometric information on the same manifold, including regions of small
measure. A uniform route through the Ricci comparison would suffice.
On a homotopy four-sphere the integral bound \eqref{eq:ricci-budget}
would also suffice. It needs no spectral input, it requires
\(\Y>8\sqrt3\,\pi\), and it is not implied by relative pinching.
Alternatively, smooth nondegenerate convergence after
diffeomorphisms of the original manifold would pass the Yamabe
minimality and spectral equalities to a limit and give
Theorem~\ref{thm:spectral}. Establishing such control for a
constructed family is a separate problem from producing trial
Rayleigh quotients.

\paragraph{Spinorial curvature and zeros.}
The canonical boundary fields give data on every input,
with a fixed end value and a smooth compact transverse pole jet.
Their interior zeros are unaffected by changing the unit source at
a fixed pole. Varying the pole or metric to control these zeros is
open. A first-order smallness argument must also supply
a derivative bound such as \eqref{eq:weyl-defect}, or produce a
nondegenerate limit of the kind in
Section~\ref{sec:twistor-variation}. The explicit oscillatory
family in Proposition~\ref{prop:oscillation} and the fixed-end
scaling of the Dirichlet energy give two independent obstructions
to deriving such control from the residual size alone. The
jump correction of Proposition~\ref{prop:dirac-jumps} can be used
once the actual seams and traces are supplied. To exploit it in the
geometric problem one must obtain compatible end identifications
and bound the projected derivative on the whole residual. The cover
estimate controls it only with a bound on the summed bulk first
derivatives and the energy of the uncovered part. The cutoff
criterion \eqref{eq:dirac-volume-witness} applies on actual marked
faces. Producing those comparisons on an approximately critical
sequence remains open.

\paragraph{Degree and the remaining eigenspaces.}
For a supplied minimizing metric, recognition from a near-zero
sequence can use retention of nonzero degree, a uniform covariance
bound in five directions, or a stress condition as in \eqref{eq:stress}.
The analytic degree cap in Theorem~\ref{thm:degree} identifies
one possible sufficient estimate. At scalar equality, the
existence or exclusion of a first eigenmap into a great
\(S^2\) or \(S^3\) on a nonround homotopy sphere remains unresolved.
The Hopf and Ricci identities delimit these cases. Even their
exclusion would leave the construction of a minimizing metric
and an eigenmap on every \(M\) to be established.

\bibliographystyle{unsrt}
\bibliography{references}

\end{document}
